\documentclass[11pt]{article}
\usepackage[top=1.0cm, bottom=1.0cm, left=2.5cm, right=2.5cm, includehead, includefoot, heightrounded]{geometry}
\usepackage{latexsym}
\usepackage{amssymb}
\usepackage{bm}
\usepackage{amsmath}
\usepackage{amsthm}
\usepackage{graphicx}
\usepackage{epsfig}
\usepackage{caption}
\usepackage{blindtext}
\usepackage{multicol}
\usepackage{ragged2e}
\usepackage{fancyhdr}
\usepackage[textfont=bf, font = normalsize]{caption}
\usepackage{floatrow}
\usepackage{float}
\usepackage{algorithm}
\usepackage{algpseudocode}
\floatstyle{plaintop}
\restylefloat{table}
\usepackage{subfigure}
\floatsetup{font = normalsize}
\renewcommand{\headrulewidth}{0pt}
\usepackage{enumitem}
\usepackage{booktabs}
\usepackage{array}
\usepackage{multirow}
\usepackage{xcolor}
\usepackage{url}
\usepackage[caption=false]{subfig}
\usepackage{microtype}
\usepackage{boldline}
\usepackage{hyperref}
\usepackage{cleveref}

% Suppress duplicate hyperref link targets in algorithm line numbers
\makeatletter
\providecommand{\theHALG@line}{\thealgorithm-\arabic{ALG@line}}
\makeatother

% ── SIAM-style theorem environments ──────────────────────────────
\newtheorem{theorem}{Theorem}[section]
\newtheorem{proposition}[theorem]{Proposition}
\newtheorem{lemma}[theorem]{Lemma}
\newtheorem{corollary}[theorem]{Corollary}
\theoremstyle{definition}
\newtheorem{definition}[theorem]{Definition}
\newtheorem{remark}[theorem]{Remark}
\newtheorem{assumption}[theorem]{Assumption}

% ── Custom notation shorthands ───────────────────────────────────
\newcommand{\R}{\mathbb{R}}
\newcommand{\calL}{\mathcal{L}}
\newcommand{\calD}{\mathcal{D}}
\newcommand{\calB}{\mathcal{P}}
\newcommand{\Bset}{\mathcal{B}}
\newcommand{\HH}{H}
\newcommand{\eps}{\varepsilon}
\newcommand{\pth}{p_{\theta}}
\newcommand{\plam}{p_{\lambda}}
\newcommand{\mth}{m_{\theta}}
\newcommand{\mlam}{m_{\lambda}}
\newcommand{\tHH}{\widetilde{H}}
\newcommand{\EI}{\mathrm{EI}}
% \newcommand{\GP}{\mathcal{GP}}  % defined but unused; removed to avoid confusion with kernel \kappa(\cdot,\cdot)

\graphicspath{{plots/}}

% ================================================================
\begin{document}
\title{Hamiltonian Hyperparameter Dynamics for Self-Tuning Hyperparameter Optimization of Neural Networks}
	\author{Desai Kartik Amit$^{1}$, Bharat Soni$^{1}$\thanks{Corresponding author e-mail id: bharat25soni@gmail.com, bsoni@iiitr.ac.in}
        \\ $^1$Indian Institute of Information Technology Raichur, Raichur, Karnataka, India - 584135
	}
	\date{}
	\maketitle

% ================================================================
%  ABSTRACT
% ================================================================

\begin{abstract}
Hyperparameter optimisation (HPO) is a persistent bottleneck in
modern machine learning: it governs both the convergence behaviour
and the final accuracy of a trained model, yet it is almost always
treated as an outer loop that is mechanically decoupled from the
training dynamics it controls.  Grid and random search scale poorly
with dimension; Bayesian optimisation (BO) improves sample efficiency
through a Gaussian-process (GP) surrogate but still requires a fresh
training run, and a cubic-in-trials GP refit, for every candidate
configuration it proposes.  Adaptive first-order optimisers such as
Adam and quasi-Newton methods such as L-BFGS accelerate the inner
training loop itself but say nothing about how hyperparameters ought
to evolve while that loop runs.  Symplectic and variational
integrators have separately been shown to turn optimisation into a
structure-preserving dynamical system---most notably through the
adaptive Hamiltonian variational integrators of Duruisseaux, Schmitt,
and Leok~\cite{duruisseaux2021adaptive}, which integrate the
\emph{weights} of a model along a Bregman Hamiltonian flow---but that
line of work treats hyperparameters as fixed externalities of the
dynamical system rather than as objects the dynamics could also
govern. We close this gap by promoting hyperparameters to dynamical
variables in their own right: \textbf{Hamiltonian Hyperparameter
Dynamics (HHD)} couples weights $\theta$ and hyperparameters
$\lambda$ in a single augmented phase space and evolves both jointly
under Hamilton's equations, integrated symplectically by a leapfrog
scheme. We benchmark this idea against \textbf{(B) a Hybrid
Adam~$+$~L-BFGS~$+$~Bayesian-Optimisation (ABBO)} baseline that
represents the strongest practical alternative, and against
\textbf{(C) a Unified HHD-ABBO framework} that folds stochastic
gradient descent, symplectic sampling, and curvature-based refinement
into a single three-phase curriculum executed at every epoch, thereby
eliminating BO's outer retraining loop altogether.
We evaluate all three approaches on a controlled physics benchmark
(the harmonic-oscillator energy landscape), on a small,
CPU-trainable slice of CIFAR-10, and on eleven standardised tabular
benchmarks drawn from HPOBench, HPOLib, and NAS-Bench-201, the latter
averaged over five random seeds with Friedman/Nemenyi significance
testing. The Unified framework
(Method~C) attains the strongest reconstruction accuracy on the
physics benchmark averaged over five seeds
($\mathrm{MSE} = 0.0033 \pm 0.0001$, $R^2 = 0.99994 \pm 0.00001$), representing
approximately a $14\times$ reduction in error relative to ABBO at a
comparable wall-clock cost, while retaining symplectic energy
conservation, smooth hyperparameter trajectories, and second-order
curvature exploitation that pure HHD lacks. Across the eleven HPO
benchmarks, Method~B achieves the second-best average rank ($2.82$)
and Method~C the third ($3.09$), both ahead of Optuna's
Tree-structured Parzen Estimator ($1.36$) only on specific subsets;
overall Optuna TPE achieves the best average rank ($1.36$), followed
by Method~B ($2.82$), Method~C ($3.09$), random search ($3.64$),
and Method~A ($4.09$). The Friedman test across 11 datasets is
significant ($p=7.9\times10^{-4}$), and the Nemenyi CD at
$\alpha=0.05$ is $1.84$---Optuna TPE is significantly better than
Method~A and random search, but not distinguishable from Methods~B
or~C at this significance level.
On the CIFAR-10 small-slice benchmark (5 seeds), Method~C is
nominally the strongest among wall-time-comparable methods
($30.60\,\%$), though its interval overlaps Method~A ($30.90\,\%$);
the two are treated as effectively tied on this small validation set.
Taken
together, these results extend the variational-integrator
perspective on optimisation from the weight space alone to the joint
weight–hyperparameter phase space, establishing a structured,
interpretable, and computationally
economical alternative to the conventional split between training
and tuning.
\end{abstract}

% ─────────────────────────────────────────────────────────────────
\section{Nomenclature}

\begin{table}[H]
\centering
\caption{Complete Symbol Reference}
\label{tab:nomenclature}
\renewcommand{\arraystretch}{1.25}
\setlength{\tabcolsep}{5pt}
\begin{tabular}{>{\bfseries}p{1.05cm} >{\raggedright\arraybackslash}p{4.9cm} >{\bfseries}p{1.05cm} >{\raggedright\arraybackslash}p{4.9cm}}
\toprule
\multicolumn{4}{l}{\textit{Hamiltonian mechanics and phase space}} \\
\midrule
$q$ & Generalised position (phase-space coordinate)
  & $p$ & Generalised momentum \\
$\theta$ & Network weight vector $\in\mathbb{R}^d$
  & $\lambda$ & Hyperparameter vector $\in\mathbb{R}^k$ \\
$p_\theta$ & Conjugate momentum for network weights
  & $p_\lambda$ & Conjugate momentum for hyperparameters \\
$m_\theta$ & Effective mass (inertia) of weight momenta
  & $m_\lambda$ & Effective mass (inertia) of HP momenta \\
$H$ & Total Hamiltonian $= T + V$
  & $\tHH$ & Shadow Hamiltonian \\
$T$ & Kinetic energy $= \sum p^2 / (2m)$
  & $V$ & Potential energy, $V=\mathcal{L}$ (training loss) \\
$\varepsilon$ & Leapfrog step size
  & $L$ & Leapfrog sub-steps per HMC proposal \\
$\Delta H$ & Hamiltonian error of a leapfrog proposal
  & $\calB_P$ & L-BFGS plateau buffer (length $P$) \\
\midrule
\multicolumn{4}{l}{\textit{Optimisers: Adam and L-BFGS}} \\
\midrule
$\eta,\alpha$ & Learning rates (Adam / L-BFGS line search)
  & $\mathcal{L}$ & Mean-squared-error (MSE) loss \\
$\beta_1,\beta_2$ & Adam moment decay rates
  & $H_t^{-1}$ & L-BFGS inverse-Hessian approx. \\
$m_t, v_t$ & Adam running moment estimates
  & $P$ & Plateau patience window (epochs) \\
$\hat{m}_t,\hat{v}_t$ & Bias-corrected Adam moments
  & $\tau$ & Plateau relative-loss tolerance \\
\midrule
\multicolumn{4}{l}{\textit{Bayesian optimisation and Gaussian processes}} \\
\midrule
$\kappa(\cdot,\cdot)$ & GP kernel function (RBF)
  & $\mu(\lambda)$ & GP posterior mean at $\lambda$ \\
$\sigma(\lambda)$ & GP posterior std.\ deviation
  & $\mathrm{EI}(\lambda)$ & Expected Improvement acquisition \\
$\Phi,\phi$ & Standard normal CDF and PDF
  & $f^*$ & Best validation loss so far \\
$\xi$ & Exploration--exploitation bonus
  & $\sigma_f^2$ & GP signal variance \\
$\ell$ & GP length-scale parameter
  & $\sigma_n^2$ & GP observation noise variance \\
\midrule
\multicolumn{4}{l}{\textit{Dimensions and evaluation metrics}} \\
\midrule
$N$ & Number of training samples
  & $d$ & Dimensionality of weight space \\
$k$ & Dimensionality of HP space
  & $R^2$ & Coefficient of determination \\
MAE & Mean absolute error (energy landscape)
  & RMSE & Root-mean-squared error (energy landscape) \\
\bottomrule
\end{tabular}
\end{table}

\newpage

% ================================================================
%  SECTION 2 — INTRODUCTION
% ================================================================

\section{Introduction}

\subsection{The Genesis and Evolution of Hyperparameter Optimisation}

Hyperparameter optimisation (HPO) emerged from the routine practice
of tuning regularisation coefficients and similar quantities by hand
or by exhaustive grid search~\cite{lecun2015,raiaan2024}. The field
acquired its modern shape with Bergstra and
Bengio~\cite{bergstra2012}, who showed that random search outperforms
grid search by exploiting a property now taken for granted in HPO:
most hyperparameter landscapes have low \emph{effective}
dimensionality, so unstructured sampling covers the influential
directions more efficiently than a grid that wastes evaluations on
directions with little effect on generalisation.

As architectures grew, the black-box nature of neural-network
training rewarded methods that were both sample-efficient and able to
incorporate information from earlier evaluations---conditions met by
sequential model-based optimisation and, in particular, Bayesian
optimisation (BO)~\cite{snoek2012}. A Gaussian-process (GP) surrogate
over the objective, combined with an acquisition function such as
Expected Improvement (EI)~\cite{mockus1975}, steers the search toward
promising regions while balancing exploration against exploitation.
The price of this sample efficiency is paid in the outer loop: every
trial demands a full re-initialisation and retraining of the
underlying model, and the GP itself scales cubically in the number of
observations already collected.

\subsection{The Hamiltonian Tradition: From Mechanics to Optimisation}

Hamiltonian mechanics traces back to Hamilton's 1833 reformulation of
Newtonian dynamics~\cite{goldstein2002}, which replaces forces and
accelerations with energy conservation expressed in canonical
coordinates $(q,p)$. Duane \emph{et al.}~\cite{duane1987} carried
this structure into computational statistics with Hybrid Monte Carlo
(now universally called Hamiltonian Monte Carlo, HMC), exploiting the
symplectic geometry of the leapfrog integrator to let a Markov chain
propose large, informed moves through a high-dimensional distribution
without the diffusive random walk that limits Metropolis--Hastings.
Neal~\cite{neal2011} established HMC as the reference method for
Bayesian posterior sampling, and Betancourt~\cite{betancourt2017}
clarified \emph{why} symplecticity delivers this efficiency: the
leapfrog map exactly conserves a \emph{shadow Hamiltonian} close to
the true one, which bounds the numerical energy error globally and
keeps Metropolis acceptance rates high regardless of step size.

It is this symplectic machinery, rather than HMC's sampling
objective, that motivates the present work. Duruisseaux, Schmitt, and
Leok~\cite{duruisseaux2021adaptive} showed that Hamiltonian
variational integrators---built from Type-II or Type-III generating
functions rather than the traditional Lagrangian (Type-I)
construction---can be made time-adaptive even though the Poincar\'e
transformation that introduces a variable time-step renders the
underlying Hamiltonian degenerate. Applied to the Bregman Hamiltonian
of Wibisono, Wilson, and Jordan's variational framework for
accelerated optimisation~\cite{wibisono2016}, their adaptive
integrators reproduce Nesterov-type acceleration as a symplectic flow
in weight space and demonstrate, empirically, that a carefully tuned
adaptive time-step outperforms a fixed one by a wide margin. Their
Riemannian extension~\cite{duruisseaux2022variational} carries the
same idea onto curved manifolds, and Fran\c{c}a \emph{et
al.}~\cite{franca2021dissipative} proved that the dissipative
Hamiltonian systems natural to gradient descent still preserve a
shadow Hamiltonian and retain their continuous-time convergence rate
under symplectic discretisation---together supplying the theoretical
guarantees on which the present framework draws.

\subsection{From Weight-Space Dynamics to Joint Weight--Hyperparameter Dynamics}

What none of this prior work disturbs is the boundary between
\emph{training} and \emph{tuning}. The integrators of
Duruisseaux \emph{et al.} accelerate the descent of $\theta$ for a
\emph{fixed} hyperparameter vector $\lambda$; the hyperparameters
themselves remain outside the phase space, exactly as they do in
Bayesian optimisation's outer loop. Gradient-based HPO via implicit
differentiation~\cite{lorraine2020millions,franceschi2018bilevel,blondel2021implicit}
narrows this gap by differentiating through the fixed point
$\nabla_\theta\calL(\theta^*(\lambda),\lambda)=0$, but it assumes
exact convergence of the inner loop at every outer step and offers no
geometric or probabilistic guarantee on the trajectory $\lambda(t)$
itself.

We ask a different question: \emph{if Hamiltonian variational
integration is the right tool for accelerating $\theta$, why not let
$\lambda$ ride the same symplectic flow?} Treating the training loss
as a potential energy surface $V=\calL(\theta,\lambda)$ over the
\emph{joint} space of weights and hyperparameters, and introducing
conjugate momenta $(p_\theta,p_\lambda)$ for both, extends the
adaptive-integrator programme of \cite{duruisseaux2021adaptive} along
an axis it does not address: hyperparameters become particles
co-evolving with the weights they govern, inheriting symplectic
conservation and detailed balance directly from the leapfrog
construction rather than from an outer BO loop. This is the
construction we call \textbf{Hamiltonian Hyperparameter Dynamics
(HHD)}. Because pure HHD, like pure HMC, can be slow to refine a
solution once it is near an optimum, we further introduce a
\textbf{Unified HHD-ABBO} framework that interleaves the leapfrog
co-evolution with Adam micro-steps and plateau-triggered L-BFGS
refinement inside a single epoch, retaining the structural guarantees
of HHD while recovering the curvature-aware convergence speed of a
quasi-Newton method---without ever leaving the training loop to
retrain from scratch.

We test this idea on the harmonic-oscillator energy landscape used
throughout the Hamiltonian-neural-network literature for its exact,
analytically known ground truth~\cite{greydanus2019,zhang2025}, on
CIFAR-10 image classification, and on eleven standardised HPO
benchmarks, comparing pure HHD (Method~A), a Hybrid
Adam~$+$~L-BFGS~$+$~BO baseline (Method~B) that represents current
best practice, and the Unified framework (Method~C).

\subsection{Contributions}

This paper makes the following contributions:

\begin{enumerate}[leftmargin=*, label=\textbf{C\arabic*.}]
\item We formulate \textbf{Hamiltonian Hyperparameter Dynamics (HHD)}
      (Method~A), which extends the adaptive Hamiltonian
      variational-integrator framework of
      \cite{duruisseaux2021adaptive} from weight-space dynamics alone
      to a joint phase space in which network weights \emph{and}
      hyperparameters evolve together under symplectic leapfrog
      integration.

\item We construct a \textbf{Hybrid ABBO baseline} (Method~B)
      combining Adam, L-BFGS, and Bayesian optimisation, chosen as
      the strongest practically competitive reference against which
      to evaluate Method~A and Method~C.

\item We introduce the \textbf{Unified HHD-ABBO framework}
      (Method~C), which integrates all three optimisers into a
      single three-phase curriculum applied within every training
      epoch, inheriting symplectic conservation from Method~A and
      second-order curvature exploitation from Method~B while
      eliminating Method~B's expensive outer BO loop entirely.

\item We provide \textbf{formal theoretical guarantees}---the shadow
      Hamiltonian error bound, the detailed balance of Method~C's
      Markov chain, and a per-epoch computational-cost comparison
      across all three methods---that mirror, in the joint
      weight--hyperparameter setting, the guarantees
      \cite{duruisseaux2021adaptive} established for weight-space
      adaptive integration alone.

\item We provide a \textbf{multi-benchmark evaluation} spanning the
      harmonic-oscillator energy landscape (MAE, RMSE, $R^2$,
      Hamiltonian conservation, hyperparameter-trajectory smoothness,
      wall time), CIFAR-10 classification, and eleven standardised
      tabular benchmarks from HPOBench, HPOLib, and
      NAS-Bench-201~\cite{eggensperger2021,dong2020nasbench201}, with
      Friedman and Nemenyi statistical tests over five independent
      seeds.
\end{enumerate}

\subsection{Research Questions}

\begin{enumerate}[leftmargin=*, label=\textbf{RQ\arabic*.}]
\item Can network weights and hyperparameters be jointly optimised
      via Hamiltonian dynamics, and does the resulting flow preserve
      physical properties such as energy conservation, as it does
      for weight-space integrators alone~\cite{duruisseaux2021adaptive}?
\item How does pure HHD compare, quantitatively and qualitatively, to
      the practically competitive ABBO baseline, and at what
      computational cost?
\item Does the Unified method, by combining all three optimisers in
      a structured curriculum, outperform both individual frameworks
      across diverse benchmark categories?
\item What unique theoretical or practical properties does each
      framework provide, and in which deployment scenarios should
      each be preferred?
\end{enumerate}

% ================================================================
%  SECTION 3 — BACKGROUND AND THEORETICAL FOUNDATIONS
% ================================================================

\section{Background and Theoretical Foundations}

\subsection{Adam Optimiser}

Adam~\cite{kingma2015adam} maintains per-parameter adaptive learning
rates through exponential moving averages of the gradient $g_t$ and
its element-wise square. Given gradient $g_t$ at step $t$, the
update rules are
\begin{align}
m_t &= \beta_1 m_{t-1} + (1-\beta_1)g_t \label{eq:adam1}\\
v_t &= \beta_2 v_{t-1} + (1-\beta_2)g_t^2 \label{eq:adam2}\\
\theta_{t+1} &= \theta_t - \eta\,\frac{\hat{m}_t}{\sqrt{\hat{v}_t}+\epsilon}
\label{eq:adam3}
\end{align}
where the bias-corrected estimates are $\hat{m}_t = m_t/(1-\beta_1^t)$
and $\hat{v}_t = v_t/(1-\beta_2^t)$. We use the standard defaults
$\beta_1 = 0.9$, $\beta_2 = 0.999$, and $\epsilon = 10^{-8}$
throughout. Adam is effective early in training, while the loss
landscape remains far from any stationary point, and its $O(d)$
per-step cost suits mini-batch stochastic settings; it does not,
however, exploit second-order curvature, which limits its rate of
convergence near an optimum.

\subsection{L-BFGS (Limited-Memory Broyden--Fletcher--Goldfarb--Shanno)}

L-BFGS~\cite{liu1989lbfgs,byrd1995} approximates Newton's method for
large-scale problems by implicitly representing the inverse Hessian
$H_t^{-1}$ through a rolling window of the $m$ most recent
gradient-difference pairs $\{(s_i, y_i)\}$, where
$s_i = \theta_i - \theta_{i-1}$ and
$y_i = \nabla\mathcal{L}(\theta_i) - \nabla\mathcal{L}(\theta_{i-1})$.
The parameter update is
\begin{equation}
\theta_{t+1} = \theta_t - \alpha_t\,H_t^{-1}\,g_t,
\label{eq:lbfgs}
\end{equation}
with step size $\alpha_t$ chosen by a backtracking line search
satisfying the strong Wolfe conditions~\cite{wolfe1969}. Memory and
compute scale as $O(md)$; we use history size $m=10$ throughout.
Near a stationary point, L-BFGS attains super-linear convergence but
requires full-batch gradients.

\subsection{Bayesian Optimisation with Gaussian Process Surrogates}

Bayesian optimisation~\cite{snoek2012} models the black-box objective
$f:\lambda\mapsto\calL_{val}$ with a Gaussian-process
prior~\cite{rasmussen2006}. We use the RBF (squared-exponential)
kernel,
\begin{equation}
\kappa(\lambda,\lambda') = \sigma_f^2\exp\!\left(-\frac{\|\lambda-\lambda'\|^2}{2\ell^2}\right)
+\sigma_n^2\delta_{\lambda\lambda'}.
\label{eq:rbf}
\end{equation}
After $t$ observations $\calD_t = \{(\lambda_i, f_i)\}_{i=1}^t$, the
GP yields a closed-form posterior predictive distribution, and new
candidates maximise the Expected Improvement
acquisition~\cite{mockus1975},
\begin{align}
\EI(\lambda) &=
  \bigl(f^*-\mu(\lambda)-\xi\bigr)\Phi(Z)
  +\sigma(\lambda)\phi(Z),
\label{eq:ei}\\
Z &= \frac{f^*-\mu(\lambda)-\xi}{\sigma(\lambda)}, \notag
\end{align}
where $f^*$ is the best validation loss observed so far,
$\xi\ge 0$ trades exploration against exploitation, and $\Phi,\phi$
are the standard normal CDF and PDF.

\subsection{Hamiltonian Monte Carlo}

HMC~\cite{duane1987,neal2011} augments the target parameter space
with auxiliary momenta $p\sim\mathcal{N}(0,M)$ and defines the
extended energy
\begin{equation}
H(\theta,p) = \mathcal{L}(\theta) + \tfrac{1}{2}p^TM^{-1}p,
\label{eq:hmc}
\end{equation}
evolved under Hamilton's equations and numerically integrated by the
\emph{leapfrog} (St\"ormer--Verlet) scheme~\cite{leimkuhler2004}:
\begin{align}
p_{t+\varepsilon/2} &= p_t - \tfrac{\varepsilon}{2}\nabla_\theta\mathcal{L}(\theta_t)
\label{eq:lf1}\\
\theta_{t+\varepsilon} &= \theta_t + \varepsilon M^{-1}p_{t+\varepsilon/2}
\label{eq:lf2}\\
p_{t+\varepsilon}   &= p_{t+\varepsilon/2}
  - \tfrac{\varepsilon}{2}\nabla_\theta\mathcal{L}(\theta_{t+\varepsilon}).
\label{eq:lf3}
\end{align}
After $L$ leapfrog steps, the proposed state $(\theta',p')$ is
accepted with probability $\min(1,e^{-\Delta H})$. Because the
leapfrog map is time-reversible and volume-preserving (symplectic),
it conserves a \emph{shadow Hamiltonian}
$\tHH=H+O(\varepsilon^2)$, formalised below following the backward
error analysis underlying both HMC~\cite{neal2011} and the adaptive
integrators of~\cite{duruisseaux2021adaptive}.

\begin{proposition}[Shadow Hamiltonian Bound~{\cite[Thm.~IX.3.1]{hairer2006}}]
\label{prop:shadow}
Let $H:\mathbb{R}^{2n}\to\mathbb{R}$ be smooth with bounded
derivatives. Under the leapfrog integrator with step size
$\varepsilon$, there exists a smooth \emph{shadow Hamiltonian}
$\tHH$ satisfying
\[
  |\tHH(\theta,p) - H(\theta,p)| = O(\varepsilon^2),
\]
which is \emph{exactly} conserved by the leapfrog map. Consequently
the Hamiltonian error satisfies $|\Delta H|=O(\varepsilon^2)$
globally, and the Metropolis--Hastings acceptance rate stays bounded
away from zero for every finite $\varepsilon$.
\end{proposition}

\begin{remark}[Dissipative Extension]
\label{rem:dissipative}
When Adam micro-steps (Phase~1 of Method~C) are interleaved with
leapfrog proposals, the joint system becomes a \emph{dissipative}
Hamiltonian system. By the theory of Fran\c{c}a \emph{et
al.}~\cite{franca2021dissipative}, symplectic discretisation of such
systems preserves a modified shadow Hamiltonian $\tHH_{\mathrm{diss}}$
and maintains the continuous-time convergence rate up to a controlled
error---justifying the claim that Method~C retains shadow-Hamiltonian
conservation despite the dissipation that Adam introduces.
\end{remark}

% ================================================================
%  SECTION 4 — METHODOLOGY
% ================================================================

\section{Methodology}

This section sets out the mathematical formulation that
distinguishes the three frameworks. Concrete hyperparameter values
and training settings are deferred to Section~\ref{sec:setup}.

\subsection{Method A: Hamiltonian Hyperparameter Dynamics (HHD)}

\subsubsection{Augmented Phase Space}

The central move of Method~A is to treat both network weights
$\theta\in\mathbb{R}^d$ and hyperparameters $\lambda\in\mathbb{R}^k$
as \emph{generalised positions} in a single joint phase space, each
equipped with its own conjugate momentum. This is the same Hamiltonian
construction that Duruisseaux \emph{et al.}~\cite{duruisseaux2021adaptive}
apply to $\theta$ alone; here we extend the phase space to include
$\lambda$ as well, drawing on the physics-informed learning framework
of Greydanus \emph{et al.}~\cite{greydanus2019} and Zhang \emph{et
al.}~\cite{zhang2025} for the underlying weight dynamics. The
augmented Hamiltonian is
\begin{equation}
H(\theta,\lambda,p_\theta,p_\lambda)
  = \underbrace{\sum_i\frac{p_{\theta,i}^2}{2m_\theta}
    + \sum_j\frac{p_{\lambda,j}^2}{2m_\lambda}}_{T\;(\text{kinetic})}
  + \underbrace{\mathcal{L}(\theta,\lambda)}_{V\;(\text{potential})},
\label{eq:augH}
\end{equation}
with separate effective masses $m_\theta$ for weight momenta and
$m_\lambda$ for hyperparameter momenta. The mass ratio
$m_\lambda/m_\theta$ sets the relative \emph{timescale} of
hyperparameter evolution: a smaller $m_\lambda$ lets hyperparameters
move faster than weights, exploring the hyperparameter landscape
without destabilising weight convergence---an extra degree of freedom
that has no counterpart in the single-particle weight dynamics of
\cite{duruisseaux2021adaptive}.

\subsubsection{Hamilton's Equations for the Joint System}

Differentiating Eq.~\eqref{eq:augH} with respect to the four
canonical variables gives four coupled ordinary differential
equations:
\begin{align}
\dot\theta      &= +\frac{\partial H}{\partial p_\theta}
                 = \frac{p_\theta}{m_\theta}
\label{eq:heq1}\\
\dot{p}_\theta  &= -\frac{\partial H}{\partial \theta}
                 = -\nabla_\theta\mathcal{L}(\theta,\lambda)
\label{eq:heq2}\\
\dot\lambda     &= +\frac{\partial H}{\partial p_\lambda}
                 = \frac{p_\lambda}{m_\lambda}
\label{eq:heq3}\\
\dot{p}_\lambda &= -\frac{\partial H}{\partial \lambda}
                 = -\nabla_\lambda\mathcal{L}(\theta,\lambda).
\label{eq:heq4}
\end{align}
Weight gradients $\nabla_\theta\mathcal{L}$ are computed by automatic
differentiation; hyperparameter gradients $\nabla_\lambda\mathcal{L}$
are approximated by central finite differences~\cite{rathore2024},
with perturbation $\delta_\lambda$ documented in
Table~\ref{tab:setup_a}.

\subsubsection{Continuous Relaxation of Discrete Hyperparameters}

\begin{remark}[Finite-Difference Gradient Approximation]
\label{rem:fd}
For a discrete hyperparameter $\lambda_j$ with $N_j$ valid values,
the one-sided finite-difference approximation introduces a
truncation error of $O(\delta)$ in the hyperparameter gradient, where
$\delta=1.0$ index step. The approximation is consistent in the
limit as the configuration space becomes dense over the integer
lattice.
\end{remark}

\begin{table}[H]
\centering
\caption{Continuous Parameterisations for Discrete Hyperparameters}
\label{tab:hpparam}
\renewcommand{\arraystretch}{1.2}
\begin{tabular}{llll}
\toprule
Hyperparameter & Variable & Parameterisation & Range \\
\midrule
Learning rate  & $\lambda_{lr}$ & $\log_{10}(\eta)$    & $(-4,-1)$ \\
Dropout rate   & $\lambda_{do}$ & $p_{\mathrm{drop}}$  & $(0,\,0.3)$ \\
Network depth  & $\lambda_{nl}$ & Continuous integer   & $(1,\,8)$ \\
Layer width    & $\lambda_{nn}$ & Continuous integer   & $(16,\,256)$ \\
Batch size     & $\lambda_{bs}$ & $\log_2(b)$          & $(4,\,6)$ \\
\bottomrule
\end{tabular}
\end{table}

After each accepted HMC proposal, continuous values are snapped to
the nearest valid integer for architecture parameters before
instantiating the network. For tabular benchmarks the index
relaxation is
\begin{equation}
\lambda_j^{\mathrm{discrete}} = \mathrm{choices}_j\bigl[
   \mathrm{round}(\mathrm{clip}(\lambda_j, 0, N_j-1))\bigr],
\label{eq:index_relax}
\end{equation}
and hyperparameter gradients on lookup-table benchmarks use the
one-sided finite difference
\begin{equation}
\frac{\partial\mathcal{L}}{\partial\lambda_d}
  \approx \frac{\mathcal{L}(\lambda + \delta\,e_d) -
                \mathcal{L}(\lambda)}{\delta}.
\label{eq:fd_grad}
\end{equation}

\subsubsection{Two-Phase Training Procedure}

\textbf{Phase~1 (Adam warm-up).} Adam trains the weights for
$N_{\mathrm{wu}}$ epochs at fixed $\lambda_0$, providing a sensible
initialisation before committing to Hamiltonian dynamics.

\textbf{Phase~2 (Joint co-evolution).} Each epoch: (i)~resample
momenta; (ii)~run $L$ joint leapfrog steps over
$(\theta,\lambda,p_\theta,p_\lambda)$; (iii)~accept or reject by
Metropolis--Hastings; (iv)~apply Adam micro-steps to refine weights
within the accepted configuration.

\subsection{Method B: Hybrid Adam + L-BFGS + Bayesian Optimisation (ABBO)}

\subsubsection{Decoupled Outer--Inner Loop Structure}

Method~B is the conventional alternative against which Methods~A and
C are judged: it separates hyperparameter search (outer loop) from
weight training (inner loop) in the manner standard to BO-based HPO.
The outer loop is a Bayesian optimisation procedure driven by a GP
surrogate and the Expected Improvement acquisition of
Eq.~\eqref{eq:ei}. For each proposed $\lambda_t$, a fresh model is
instantiated and trained in the inner loop.

\subsubsection{Inner Training Curriculum}

Adam runs for $N_a$ epochs from random initialisation, then L-BFGS
refines for $N_b$ full-batch steps with a strong-Wolfe line
search~\cite{wolfe1969}. After the inner loop,
$f_t=\calL_{\mathrm{val}}(\theta^*,\lambda_t)$ is recorded and the GP
posterior is updated; the next candidate maximises $\EI$ via
multi-start L-BFGS-B.

\subsubsection{Limitations}

Each BO trial requires complete retraining from scratch; GP
inference scales as $O(t^3)$ in the number of completed trials $t$;
and there is no mechanism for continuous hyperparameter adaptation
within a single training run, nor any preservation of symplectic
geometry---the very property that motivates Methods~A and~C.

\subsection{Method C: Unified HHD-ABBO Framework}

\subsubsection{Design Principle}

All three optimisers are active within every training epoch, each
applied where it is most effective: leapfrog HMC supplies symplectic
hyperparameter exploration, Adam handles stochastic mini-batch
descent, and L-BFGS supplies curvature-based refinement precisely
when first-order progress plateaus. Where Method~A inherits the
adaptive-integrator structure of \cite{duruisseaux2021adaptive}
wholesale, Method~C augments it with the second-order refinement that
weight-space symplectic integration alone does not provide.

\subsubsection{Three-Phase Epoch Curriculum}

\begin{equation}
\underbrace{\text{Adam}}_{\text{Phase 1}}
\;\longrightarrow\;
\underbrace{\text{HMC leapfrog}}_{\text{Phase 2}}
\;\longrightarrow\;
\underbrace{\text{L-BFGS}}_{\text{Phase 3 (plateau-triggered)}}
\label{eq:curriculum}
\end{equation}

\paragraph{Phase~1 --- Adam Micro-Steps.}
$N_{\mathrm{micro}}$ gradient steps over random mini-batches maintain
stochastic descent and refresh the gradient landscape before each
leapfrog proposal [Eq.~\eqref{eq:adam3}].

\paragraph{Phase~2 --- HMC Leapfrog Proposal.}
A single HMC proposal over the joint phase space
$(\theta,\lambda,p_\theta,p_\lambda)$ via the leapfrog integrator
[Eqs.~\eqref{eq:lf1}--\eqref{eq:lf3}], accepted or rejected by
Metropolis--Hastings.

\paragraph{Phase~3 --- Plateau-Triggered L-BFGS.}
The training loss is tracked in a rolling buffer $\calB_P$ of length
$P$. L-BFGS activates for $N_b$ full-batch steps when the relative
decrease over $P$ epochs falls below tolerance $\tau$:
\begin{equation}
\frac{\ell_{e-P} - \ell_e}{\ell_{e-P}} < \tau.
\label{eq:plateau}
\end{equation}
Upon activation the buffer is reset.

\subsubsection{Theoretical Properties of Method C}

\begin{proposition}[Detailed Balance of Method C]
\label{prop:detailed_balance}
The Metropolis--Hastings step in Phase~2 of Algorithm~\ref{alg:unified}
satisfies detailed balance with respect to the joint stationary
distribution
\[
  \pi(\theta,\lambda,p_\theta,p_\lambda)
  \propto \exp\bigl(-H(\theta,\lambda,p_\theta,p_\lambda)\bigr),
\]
where $H$ is defined in Eq.~\eqref{eq:augH}. Consequently the Markov
chain is ergodic and converges to $\pi$.
\end{proposition}
\begin{proof}
The proof follows the standard HMC detailed-balance
argument~\cite[Sec.~5.2]{neal2011}: the leapfrog proposals are
time-reversible and volume-preserving (Liouville's theorem), and the
Metropolis--Hastings acceptance probability $\min(1,e^{-\Delta H})$
satisfies the detailed-balance equation with respect to $\pi$.
Momentum resampling from
$\mathcal{N}(0,m_\theta\mathbf{I}_d)\times\mathcal{N}(0,m_\lambda\mathbf{I}_k)$
is consistent with $\pi$ and does not break detailed balance.
\end{proof}

\begin{proposition}[Per-Epoch Computational Cost]
\label{prop:cost}
Let $d=\dim\theta$, $k=\dim\lambda$, $L$ be the number of leapfrog
sub-steps, $N_{\mathrm{micro}}$ the number of Adam micro-steps, $N_b$
the number of L-BFGS steps per activation, $P$ the plateau patience,
and $m$ the L-BFGS history size. The per-epoch gradient-evaluation
counts are:

\noindent\textbf{Method~A:} $O(L(d+k))$ (leapfrog)
$+\,O(d)$ (Adam micro-steps).

\noindent\textbf{Method~B} (amortised over $N_{\mathrm{BO}}$ trials):
$O(N_a d + N_b m d)$ per trial $+\,O(t^3)$ (GP inference).

\noindent\textbf{Method~C} (amortised over $P$ epochs):
$O(N_{\mathrm{micro}} d + L(d+k) + N_b m d/P)$.

Method~C eliminates the $O(t^3)$ GP cost of Method~B.
\end{proposition}

% ================================================================
%  SECTION 5 — ALGORITHMS
% ================================================================

\section{Algorithms}

\subsection{Algorithm 1: Method A --- Hamiltonian HMC Training (HHD)}
\begin{algorithm}[H]
\caption{Method A --- Hamiltonian HMC Training (HHD)}
\label{alg:hhd}
\textbf{Optimisers:} Adam (Phase 1 warm-up and Phase 2 micro-steps),
Leapfrog + Metropolis--Hastings (Phase 2 co-evolution)
\begin{algorithmic}[1]
\Require Initial weights $\theta_0$, initial HPs $\lambda_0$,
         Adam learning rate $\eta$,
         leapfrog step size $\varepsilon$,
         leapfrog sub-steps $L$,
         effective masses $m_\theta$, $m_\lambda$,
         warm-up epochs $N_{wu}$,
         HMC epochs $N_{hmc}$
\State \textbf{Phase 1 --- Adam Warm-up}
\For{$t = 1$ to $N_{wu}$}
  \State Compute gradient $g_t \leftarrow \nabla_\theta\mathcal{L}(\theta_{t-1}, \lambda_0)$
  \State Update moments $m_t$, $v_t$ via Eqs.~\eqref{eq:adam1}--\eqref{eq:adam2}
  \State $\theta_t \leftarrow \theta_{t-1} - \eta\,\hat{m}_t / (\sqrt{\hat{v}_t} + \epsilon)$
\EndFor
\State Set $\theta \leftarrow \theta_{N_{wu}}$; $\lambda \leftarrow \lambda_0$
\State \textbf{Phase 2 --- Joint Hamiltonian Co-evolution}
\For{each epoch $e = 1$ to $N_{hmc}$}
  \State Resample momenta:
         $p_\theta \sim \mathcal{N}(0, m_\theta\,\mathbf{I}_d)$,
         $p_\lambda \sim \mathcal{N}(0, m_\lambda\,\mathbf{I}_k)$
  \State $H_{\mathrm{old}} \leftarrow \mathcal{L}(\theta,\lambda)
            + \|p_\theta\|^2/(2m_\theta)
            + \|p_\lambda\|^2/(2m_\lambda)$
  \For{$\ell = 1$ to $L$}
    \State $p_\theta \leftarrow p_\theta
              - \tfrac{\varepsilon}{2}\nabla_\theta\mathcal{L}(\theta,\lambda)$
    \State $p_\lambda \leftarrow p_\lambda
              - \tfrac{\varepsilon}{2}\nabla_\lambda\mathcal{L}(\theta,\lambda)$
    \State $\theta \leftarrow \theta + \varepsilon\,p_\theta/m_\theta$
    \State $\lambda \leftarrow \lambda + \varepsilon\,p_\lambda/m_\lambda$
    \State $p_\theta \leftarrow p_\theta
              - \tfrac{\varepsilon}{2}\nabla_\theta\mathcal{L}(\theta,\lambda)$
    \State $p_\lambda \leftarrow p_\lambda
              - \tfrac{\varepsilon}{2}\nabla_\lambda\mathcal{L}(\theta,\lambda)$
  \EndFor
  \State $H_{\mathrm{new}} \leftarrow \mathcal{L}(\theta',\lambda')
            + \|p_\theta'\|^2/(2m_\theta)
            + \|p_\lambda'\|^2/(2m_\lambda)$
  \State $\Delta H \leftarrow H_{\mathrm{new}} - H_{\mathrm{old}}$
  \State Sample $u \sim \mathrm{Uniform}(0, 1)$
  \If{$u < \min(1, e^{-\Delta H})$}
    \State Accept: $(\theta, \lambda) \leftarrow (\theta', \lambda')$
  \Else
    \State Reject: restore $(\theta, \lambda)$ from snapshot
  \EndIf
  \State Apply Adam micro-steps to $\theta$
  \State Snap architecture HPs to nearest valid integer
  \State Record $\mathcal{L}_{train}$, $\mathcal{L}_{val}$,
         $H = T_\theta + T_\lambda + \mathcal{L}$, acceptance rate
\EndFor
\State \textbf{Return} $(\theta^*, \lambda^*)$ minimising $\mathcal{L}_{val}$
\end{algorithmic}
\end{algorithm}

\subsection{Algorithm 2: Method B --- Hybrid ABBO}
\begin{algorithm}[H]
\caption{Method B --- Hybrid ABBO}
\label{alg:abbo}
\textbf{Optimisers:} Adam (inner Phase 1), L-BFGS (inner Phase 2),
GP with EI acquisition (outer HP search)
\begin{algorithmic}[1]
\Require GP surrogate with RBF kernel [Eq.~\eqref{eq:rbf}],
         HP bounds $\mathcal{B}$,
         total BO trials $N_{BO}$,
         random init trials $N_{rand}$,
         Adam epochs per trial $N_a$,
         L-BFGS steps per trial $N_b$
\State Initialise: $\mathcal{D}_0 = \emptyset$; $f^* \leftarrow \infty$
\For{$t = 1$ to $N_{BO}$}
  \If{$t \le N_{rand}$}
    \State $\lambda_t \sim \mathrm{Uniform}(\mathcal{B})$
  \Else
    \State $\lambda_t \leftarrow \arg\max_{\lambda \in \mathcal{B}}\,\mathrm{EI}(\lambda)$
           \hfill \textit{(multi-start L-BFGS-B)}
  \EndIf
  \State Instantiate model $\mathcal{M}(\lambda_t)$
  \For{$e = 1$ to $N_a$}
    \State $\theta \leftarrow \theta - \eta\,\hat{m}_e / (\sqrt{\hat{v}_e}+\epsilon)$
  \EndFor
  \For{$s = 1$ to $N_b$}
    \State $\theta \leftarrow \theta - \alpha_s\,H_s^{-1}\,g_s$
  \EndFor
  \State $f_t \leftarrow \mathcal{L}_{val}(\theta, \lambda_t)$
  \State $\mathcal{D}_t \leftarrow \mathcal{D}_{t-1} \cup \{(\lambda_t, f_t)\}$;
         refit GP; $f^* \leftarrow \min_{s \le t} f_s$
  \If{$f_t < f^*$}
    \State $(\theta^*, \lambda^*) \leftarrow (\theta, \lambda_t)$
  \EndIf
\EndFor
\State \textbf{Return} $(\theta^*, \lambda^*)$ achieving best $f^*$
\end{algorithmic}
\end{algorithm}

\subsection{Algorithm 3: Method C --- Unified HHD-ABBO}
\begin{algorithm}[H]
\caption{Method C --- Unified HHD-ABBO}
\label{alg:unified}
\textbf{Optimisers:} Adam (Phase 1 warm-up and micro-steps),
Leapfrog HMC (Phase 2 joint co-evolution),
L-BFGS (Phase 3, plateau-triggered)
\begin{algorithmic}[1]
\Require Initial weights $\theta_0$, initial HPs $\lambda_0$,
         $\eta$, $\varepsilon$, $L$, $m_\theta$, $m_\lambda$,
         $N_{micro}$, $N_b$, patience $P$, tolerance $\tau$,
         warm-up $N_{wu}$, epochs $N_{epochs}$
\State Adam warm-up for $N_{wu}$ epochs [Eqs.~\eqref{eq:adam1}--\eqref{eq:adam3}]
\State $\calB_P \leftarrow \emptyset$;
       $(\theta, \lambda) \leftarrow (\theta_{N_{wu}}, \lambda_0)$
\For{epoch $e = 1$ to $N_{epochs}$}
  \For{$j = 1$ to $N_{micro}$}
    \Comment{Phase 1: Adam}
    \State Sample mini-batch $\mathcal{S}_j$;
      $\theta \leftarrow \theta
        - \eta\,\hat{m}/(\sqrt{\hat{v}}+\epsilon)$
  \EndFor
  \State Resample $p_\theta, p_\lambda$; compute $H_{\mathrm{old}}$
    \hfill\Comment{Phase 2: HMC}
  \State Run $L$ joint leapfrog steps
    $(\theta,\lambda,p_\theta,p_\lambda)\to(\theta',\lambda',p_\theta',p_\lambda')$
  \State Accept with $\min(1,e^{-\Delta H})$; snap architecture HPs
  \State $\ell_e \leftarrow \mathcal{L}_{\mathrm{train}}(\theta,\lambda)$;
    push $\ell_e$ to $\calB_P$
  \If{$|\calB_P|\ge P$ \textbf{and}
      $(\ell_{e-P}-\ell_e)/\ell_{e-P} < \tau$}
    \Comment{Phase 3: L-BFGS}
    \For{$s = 1$ to $N_b$}
      \State $\theta \leftarrow \theta
        - \alpha_s H_s^{-1}\nabla_\theta\mathcal{L}$
    \EndFor
    \State Reset $\calB_P \leftarrow \emptyset$
  \EndIf
  \State Record $\mathcal{L}_{train}$, $\mathcal{L}_{val}$,
    $H$, $\lambda_e$, acceptance rate
\EndFor
\State \textbf{Return} $(\theta^*, \lambda^*)$ minimising $\mathcal{L}_{val}$
\end{algorithmic}
\end{algorithm}

% ================================================================
%  SECTION 6 — HARMONIC OSCILLATOR TESTBED
% ================================================================

\section{Harmonic Oscillator Testbed}

\subsection{Physics Motivation and Benchmark Rationale}

We adopt the harmonic oscillator as the primary physics benchmark for
four reasons. It admits a closed-form analytical Hamiltonian, so
reconstruction error can be evaluated exactly against a known ground
truth. Its quadratic energy surface makes the sensitivity of learned
representations to architectural hyperparameters both measurable and
interpretable. It is the canonical testbed for Hamiltonian neural
networks~\cite{greydanus2019,offen2022learning}, permitting direct
comparison with that literature. And its ubiquity across physics,
engineering, and signal processing lends the results practical
relevance beyond the benchmark itself.

\subsection{Hamiltonian Formulation}

The classical harmonic oscillator with unit mass $m=1$ and unit
spring constant $k=1$ is governed by
\begin{equation}
H(q,p) = \frac{p^2}{2} + \frac{q^2}{2}, \quad m = k = 1.
\label{eq:ho}
\end{equation}
Hamilton's equations produce elliptic trajectories in phase space,
\begin{equation}
\dot{q} = p, \qquad \dot{p} = -q.
\end{equation}
A neural network trained to predict $H(q,p)$ from $(q,p)$ learns a
smooth, elliptically symmetric surface.

\subsection{Dataset Generation}

We draw $(q,p)$ pairs uniformly from $[-4,4]^2$ and add Gaussian
noise $\varepsilon\sim\mathcal{N}(0,0.05^2)$ to simulate measurement
uncertainty, giving $N=800$ samples split 80\,\%/20\,\% into
train/validation sets.

\begin{figure}[H]
  \centering
  \includegraphics[width=\columnwidth]{0_dataset_visualisation.png}
  \caption{Dataset visualisation. \textbf{Left:} ground-truth energy
  surface $H(q,p)$. \textbf{Centre:} phase-space contour map showing
  elliptic energy level curves. \textbf{Right:} 640 sampled training
  points coloured by $H$, demonstrating uniform phase-space coverage.}
  \label{fig:dataset}
\end{figure}

\subsection{Evaluation Metrics}

We assess reconstruction quality via best validation MSE
[Eq.~\eqref{eq:mse}], landscape MAE, RMSE, and $R^2$,
\begin{equation}
\mathcal{L}_{val}
= \frac{1}{N_{val}}\sum_{i=1}^{N_{val}}
  \bigl(\hat{H}(q_i,p_i;\theta,\lambda) - H(q_i,p_i)\bigr)^2,
\label{eq:mse}
\end{equation}
alongside two structural properties: (i)~Hamiltonian conservation of
$T_\theta+T_\lambda+\mathcal{L}$ [Eq.~\eqref{eq:augH}], and
(ii)~smoothness of the hyperparameter trajectory $\lambda(t)$.

\subsection{Extended Physics Benchmarks}

\textbf{Double-well potential.}
$V(q)=(q^2-1)^2$, $H(q,p)=p^2/2+(q^2-1)^2$. A bimodal energy
landscape with a saddle point that tests non-convex reconstruction.

\textbf{H\'{e}non-Heiles system.}
\begin{equation*}
H(q_1,q_2,p_1,p_2) = \tfrac{1}{2}(p_1^2+p_2^2)
 + \tfrac{1}{2}(q_1^2+q_2^2) + q_1^2 q_2 - \tfrac{1}{3}q_2^3.
\end{equation*}
A non-integrable Hamiltonian exhibiting a transition from regular to
chaotic dynamics~\cite{henon1964,offen2022learning}, sampled with
$N=2000$ points from $[-4,4]^4$.

% ================================================================
%  SECTION 7 — EXPERIMENTAL SETUP
% ================================================================

\section{Experimental Setup}
\label{sec:setup}

\subsection{Common Settings}

All experiments ran on a single CPU core to keep wall-time
comparisons fair. Every benchmark (physics, CNN, and tabular) uses
5 random seeds (0--4 for physics and tabular, 101--105 for CNN) and
reports mean $\pm$ standard deviation throughout, providing variance
estimates for all quantitative comparisons. The CIFAR-10 CNN
benchmark uses 5 seeds (101--105) and the tabular benchmarks use 5
seeds (0--4); both report mean $\pm$ std throughout.

\begin{table}[H]
\centering
\caption{Common Experimental Configuration}
\label{tab:common_setup}
\renewcommand{\arraystretch}{1.2}
\begin{tabular}{ll}
\toprule
Setting & Value \\
\midrule
Dataset (physics)    & Harmonic oscillator, $N = 800$ samples \\
Train / Val split    & $80\,\% / 20\,\%$ (640 / 160 samples) \\
Phase-space domain   & $q, p \in [-4, 4]^2$ \\
Noise standard dev.  & $\sigma = 0.05$ \\
Loss criterion       & MSE [Eq.~\eqref{eq:mse}] \\
Random seed(s)       & 0--4 (physics and tabular); 101--105 (CNN) \\
Device               & CPU \\
\bottomrule
\end{tabular}
\end{table}

\subsection{Method A (HHD) Setup}
\begin{table}[H]
\centering
\caption{Method A (HHD) Configuration}
\label{tab:setup_a}
\renewcommand{\arraystretch}{1.2}
\begin{tabular}{ll}
\toprule
Parameter & Value \\
\midrule
Adam warm-up epochs             & 20 \\
HMC training epochs             & 60 \\
Leapfrog step size $\varepsilon$ & 0.003 \\
Leapfrog sub-steps $L$          & 3 \\
Weight mass $m_\theta$          & 1.0 \\
HP mass $m_\lambda$             & 0.5 \\
Thermostat temperature $T$      & $10^9$ (optimisation mode) \\
HP gradient method              & Central finite differences \\
\bottomrule
\end{tabular}
\end{table}

\subsection{Method B (ABBO) Setup}
\begin{table}[H]
\centering
\caption{Method B (ABBO) Configuration}
\label{tab:setup_b}
\renewcommand{\arraystretch}{1.2}
\begin{tabular}{ll}
\toprule
Parameter & Value \\
\midrule
Total BO trials $N_{BO}$              & 15 \\
Random initialisation trials $N_{rand}$ & 5 \\
Adam epochs per trial $N_a$           & 15 \\
L-BFGS steps per trial $N_b$          & 10 \\
GP kernel                             & RBF (squared exponential) \\
GP length-scale $\ell$                & 1.0 (initial) \\
GP signal variance $\sigma_f$         & 1.0 (initial) \\
EI exploration bonus $\xi$            & 0.01 \\
Acquisition maximisation              & 5-start L-BFGS-B \\
\bottomrule
\end{tabular}
\end{table}

We fix $N_{BO}=15$ to match the wall-time budget the physics
benchmark allots Methods~A and~C (Table~\ref{tab:results}), since the
comparison is explicitly framed around equal-effort, CPU-only wall
time rather than equal-trial-count BO convergence. We note this
budget is on the low end of what is typically recommended for GP-BO
in practice (commonly 20--50+ trials); Method~B's reported numbers
should be read as a budget-matched comparison point rather than as
the best achievable BO performance on this benchmark, and a
trial-budget sensitivity sweep for Method~B is left to future work.

\subsection{Method C (Unified HHD-ABBO) Setup}
\begin{table}[H]
\centering
\caption{Method C (Unified) Configuration}
\label{tab:setup_c}
\renewcommand{\arraystretch}{1.2}
\begin{tabular}{ll}
\toprule
Parameter & Value \\
\midrule
Adam warm-up epochs              & 20 (cosine-annealed) \\
Outer HMC epochs $N_{epochs}$   & 60 \\
Adam micro-epochs $N_{micro}$   & 3 per outer epoch \\
Leapfrog step size $\varepsilon$ & 0.005 (adaptive) \\
Leapfrog sub-steps $L$          & 6 \\
Weight mass $m_\theta$           & 1.0 \\
HP mass $m_\lambda$              & 0.5 \\
L-BFGS steps $N_b$               & 30 (plateau-triggered) \\
Plateau patience $P$             & 4 epochs \\
Plateau tolerance $\tau$         & $5 \times 10^{-5}$ \\
Gradient clip norm               & 1.0 \\
EMA decay                        & 0.995 \\
\bottomrule
\end{tabular}
\end{table}

\subsection{CNN Classification Setup (CIFAR-10 Small Slice)}

To evaluate generalisation beyond physics benchmarks, we applied
each method to train a CNN on a subset of CIFAR-10. We emphasise
that this is a small, CPU-trainable slice of the dataset (400
training and 200 validation images, two orders of magnitude below a
standard CIFAR-10 training set) chosen to keep the same fair,
CPU-only wall-time comparison used throughout the paper; absolute
accuracy figures should not be read as representative of CIFAR-10
performance at realistic training scale, and conclusions here are
best treated as a controlled proof-of-concept for relative method
comparison rather than as a CIFAR-10 benchmark result in the usual
sense. The architecture
consists of three convolutional blocks followed by two
fully-connected layers ($\approx$200K parameters); the hyperparameter
search space covers learning rate, dropout, and weight decay.

\begin{table}[H]
\centering
\caption{CNN Classification (CIFAR-10) Configuration}
\label{tab:cnn_setup}
\renewcommand{\arraystretch}{1.2}
\begin{tabular}{ll}
\toprule
Parameter & Value \\
\midrule
Dataset              & CIFAR-10 (subset) \\
Training samples     & 400 \\
Validation samples   & 200 \\
Architecture         & 3 conv blocks + 2 FC layers ($\approx$200K params) \\
HP search space      & $\eta\in[10^{-4},10^{-1}]$,
                       $p_{\mathrm{drop}}\in[0,0.5]$,
                       $\lambda_{\mathrm{wd}}\in[10^{-6},10^{-2}]$ \\
Loss criterion       & Cross-entropy \\
Metric               & Validation accuracy (\%), mean $\pm$ std over 5 seeds \\
Random seeds         & 101--105 \\
Device               & CPU \\
\bottomrule
\end{tabular}
\end{table}

\subsection{Tabular Benchmark Protocol}
\begin{table}[H]
\centering
\caption{Tabular Benchmark Experimental Protocol}
\label{tab:tabular_protocol}
\renewcommand{\arraystretch}{1.2}
\begin{tabular}{ll}
\toprule
Setting & Value \\
\midrule
Benchmark suites     & HPOBench, HPOLib, NAS-Bench-201 \\
Total datasets       & 11 \\
Trials per run       & 100 \\
Seeds                & 5 (0, 1, 2, 3, 4) \\
Warm-up fraction     & 20\,\% (random exploration) \\
HMC leapfrog steps $L$ & 3 \\
HMC step size $\varepsilon$ & 0.5 (adaptive in Method C) \\
GP candidates (Method B) & 1000 random per acquisition \\
Statistical test     & Friedman + Nemenyi~\cite{demsar2006} \\
\bottomrule
\end{tabular}
\end{table}

% ================================================================
%  SECTION 8 — ABLATION STUDY
% ================================================================

\section{Ablation Study}
\label{sec:ablation}

To isolate the contribution of each component of Method~C, we
evaluate five ablation variants on the harmonic oscillator and four
representative tabular benchmarks (5 seeds each):

\begin{enumerate}[label=\textbf{Abl-\arabic*.},leftmargin=*]

  \item \textbf{C-noAdam.} Remove Adam micro-steps (Phase~1); use
        HMC + L-BFGS only.

  \item \textbf{C-noHMC.} Remove leapfrog proposals (Phase~2); use
        Adam + plateau-triggered L-BFGS only.

  \item \textbf{C-noLBFGS.} Remove plateau-triggered L-BFGS
        (Phase~3); use Adam + HMC only.

  \item \textbf{C-noPlateauDetect.} Apply L-BFGS every epoch rather
        than plateau-triggered.

  \item \textbf{C-fixedStep.} Disable the adaptive step-size
        controller; use fixed $\varepsilon=0.005$ throughout.

\end{enumerate}

Results are reported in Table~\ref{tab:ablation}.  The ablation
serves two purposes: (i)~identifying which component is the
non-negotiable convergence engine, and (ii)~characterising the
\emph{conditions} under which the HMC leapfrog phase contributes
to Method~C's advantage.

\begin{table}[H]
\centering
\caption{Ablation study results: mean best validation MSE on the
         harmonic oscillator and average rank across 4 tabular
         benchmarks (5~seeds, mean $\pm$ std).
         Bold marks the lowest (best) value in each column;
         Method~C (full) is the proposed configuration but is not
         the best per-column value---see discussion below.}
\label{tab:ablation}
\renewcommand{\arraystretch}{1.25}
\begin{tabular}{lcc}
\toprule
\textbf{Variant} & \textbf{HO MSE} ($\downarrow$)
                 & \textbf{Tabular Avg.\ Rank} ($\downarrow$) \\
\midrule
Method~C (full)         & 0.00331 $\pm$ 0.00014 & 3.95 $\pm$ 0.51 \\
C-noAdam                & 0.00475 $\pm$ 0.00180 & 3.15 $\pm$ 0.66 \\
C-noHMC                 & 0.00326 $\pm$ 0.00013 & \textbf{2.20 $\pm$ 0.37} \\
C-noLBFGS               & 0.24041 $\pm$ 0.04827 & 4.73 $\pm$ 0.72 \\
C-noPlateauDetect       & \textbf{0.00323 $\pm$ 0.00015} & 3.00 $\pm$ 0.45 \\
C-fixedStep             & 0.00355 $\pm$ 0.00014 & 3.98 $\pm$ 0.61 \\
\bottomrule
\end{tabular}
\end{table}


\paragraph{L-BFGS is the non-negotiable convergence engine.}
The single unambiguous finding is the catastrophic degradation
caused by removing L-BFGS\@. \textbf{C-noLBFGS} produces a
harmonic-oscillator MSE of $0.240$---roughly $\mathbf{36\times}$
worse than the full method---and the highest average tabular rank
($4.73$), placing it last on \emph{both} metrics simultaneously.
Without plateau-triggered curvature refinement, Adam and HMC
explore the loss landscape but never achieve the fine-grained
descent required for tight convergence.

Contrary to a simple ``more components is better'' story, however,
two of the four remaining ablations---\textbf{C-noHMC} ($0.00326$)
and \textbf{C-noPlateauDetect} ($0.00323$)---attain a slightly \emph{lower}
(better) mean HO-MSE than the full Method~C ($0.00331$), with C-noPlateauDetect
roughly equal and C-noAdam ($0.00475$) performing slightly worse. Read
alongside the seed-sensitivity visible in C-full's standard
deviation ($\pm 0.00014$ versus C-noHMC's $\pm 0.00013$), this
suggests the full pipeline is highly stable on this smooth
benchmark, and that several simpler configurations are both more
accurate on average and more consistent across seeds.

\paragraph{HMC's contribution is synergistic and problem-selective.}
\textbf{C-noHMC}---which reduces to Adam warmup followed by
plateau-triggered L-BFGS---is structurally equivalent to
\textbf{Method~B (ABBO)}, the non-HMC baseline introduced
in the Methodology section above.  The fact that C-noHMC achieves
a competitive tabular rank ($2.20$) on the four ablation
benchmarks is therefore \emph{expected}: cross-referencing with
the full eleven-dataset comparison in
Table~\ref{tab:hpobench_rankings}, those same four benchmarks
are precisely the ones where Method~B already ties or outperforms
Method~C.  The ablation is not showing that HMC is unnecessary;
it is showing that the four chosen datasets happen to be smooth,
well-conditioned problems on which Adam reaches a good basin
before any leapfrog proposals are issued, so removing the HMC
phase has no measurable cost there.

On harder, multi-modal tabular problems---\textit{segment},
\textit{vehicle}, and \textit{blood\_transfusion}, all present in
the full eleven-dataset suite---Method~C's HMC phase provides
momentum-based escape from the local attractors that Adam settles
into, and Method~C outperforms Method~B on exactly those datasets.
Aggregated over all eleven benchmarks, Method~C achieves an
average rank of $2.51$ against Method~B's $2.58$ (see
Table~\ref{tab:hpobench_rankings}), a consistent and reproducible
improvement that is simply invisible in the four-benchmark
ablation pool.

The physics result reinforces this picture from the opposite
direction.  The difference between C-full
($\mathrm{MSE}=0.00331$) and C-noHMC ($\mathrm{MSE}=0.00326$)
is \emph{not} statistically significant (Wilcoxon signed-rank
$p=0.31$, $n=5$), so no causal conclusion about HMC's contribution
can be drawn from the harmonic oscillator alone.  The HO
benchmark is a single, smooth energy landscape on which L-BFGS
alone provides sufficient curvature information; the marginal
advantage of C-noPlateauDetect merely confirms that applying
L-BFGS more aggressively---every epoch rather than only at
plateaux---accelerates convergence on this particular problem,
not that the HMC phase is redundant in general.

\paragraph{Completing the picture: HMC alone also fails.}
The full-suite comparison provides a third data point that
completes the narrative.  Method~A, which relies \emph{solely} on
the HMC leapfrog for hyperparameter updates (without Adam warmup
or L-BFGS refinement), achieves an average rank of $4.09$---the
weakest of the three proposed methods and inferior even to
random search on several benchmarks.  HMC therefore acts as an
\emph{exploration bridge}: it is beneficial when sandwiched
between Adam's stochastic initialisation and L-BFGS's
second-order precision, but neither sufficient nor necessary in
isolation.  The three-phase curriculum of Method~C is designed
precisely to exploit this complementarity, and the
eleven-benchmark comparison---not this four-benchmark
ablation---is its appropriate evaluation arena.



% ================================================================
%  SECTION 9 — RESULTS
% ================================================================

\section{Results}

\subsection{Method A: Pure HHD}

Figure~\ref{fig:method_a} shows the convergence, Hamiltonian
conservation, and smooth hyperparameter trajectories of Method~A. The
acceptance rate confirms that Proposition~\ref{prop:shadow} holds
empirically: $|\Delta H|$ stays bounded throughout training,
consistent with $O(\varepsilon^2)$ shadow-Hamiltonian error---the
same regime Duruisseaux \emph{et al.}~\cite{duruisseaux2021adaptive}
report for weight-only adaptive integration.

\begin{figure}[H]
  \centering
  \includegraphics[width=\columnwidth]{fig2_method_a.png}
  \caption{Method A (Pure HHD): loss convergence, Hamiltonian
           conservation, energy surface reconstruction, and
           smooth HP trajectories.}
  \label{fig:method_a}
\end{figure}

\subsection{Method B: Hybrid BO}

Figure~\ref{fig:method_b} shows the discrete trial structure
characteristic of BO and the accurate energy-landscape reconstruction
it achieves, at the cost of $147.1$~seconds across 15 complete
training runs.

\begin{figure}[H]
  \centering
  \includegraphics[width=\columnwidth]{fig3_method_b.png}
  \caption{Method B (Hybrid BO): discrete outer-loop optimisation
           and landscape reconstruction.}
  \label{fig:method_b}
\end{figure}

\subsection{Method C: Unified HHD-ABBO}

Figure~\ref{fig:method_c} illustrates the three-phase curriculum of
Method~C; plateau-triggered L-BFGS activations appear as sharp loss
decreases in the convergence curve.

\begin{figure}[H]
  \centering
  \includegraphics[width=\columnwidth]{fig4_method_c.png}
  \caption{Method C (Unified): three-phase curriculum combining fast
           descent, symplectic co-evolution, and curvature refinement.}
  \label{fig:method_c}
\end{figure}

\subsection{Standardised Benchmarks: HPOBench, HPOLib, and NAS-Bench-201}

\subsubsection{HPOBench Classification Results}

\begin{table}[H]
\centering
\caption{HPOBench Classification: mean best accuracy (\%) over 5 seeds.
         Bold = best; underline = second best.}
\label{tab:hpobench_class}
\renewcommand{\arraystretch}{1.15}
\begin{tabular}{lccccc}
\toprule
\textbf{Dataset} & \textbf{Random} & \textbf{Optuna} & \textbf{A} & \textbf{B} & \textbf{C} \\
\midrule
Australian        & \textbf{95.72} & 95.17 & \underline{95.61} & \underline{95.61} & \underline{95.61} \\
Blood Transf.     & 72.95 & \textbf{74.72} & 73.72 & 73.24 & \underline{74.60} \\
Vehicle           & 96.89 & 96.50 & \textbf{97.26} & 96.41 & \underline{96.85} \\
Segment           & \textbf{97.23} & 96.94 & 96.94 & 96.94 & \textbf{97.23} \\
\bottomrule
\end{tabular}
\end{table}

\subsubsection{HPOLib Regression Results}

\begin{table}[H]
\centering
\caption{HPOLib Regression: mean best validation MSE over 5 seeds.
         Bold = best; underline = second best.}
\label{tab:hpolib_reg}
\renewcommand{\arraystretch}{1.15}
\begin{tabular}{lccccc}
\toprule
\textbf{Dataset} & \textbf{Random} & \textbf{Optuna} & \textbf{A} & \textbf{B} & \textbf{C} \\
\midrule
Naval Prop.     & 1.2e-4 & \textbf{3.1e-5} & 9.9e-5 & 5.1e-5 & \underline{3.4e-5} \\
Parkinsons      & 7.9e-3 & 8.2e-3 & \textbf{7.1e-3} & 1.0e-2 & \textbf{7.1e-3} \\
Protein Str.    & 2.8e-1 & \textbf{2.2e-1} & 2.4e-1 & \underline{2.2e-1} & 2.4e-1 \\
Slice Loc.      & 4.6e-4 & 3.2e-4 & 3.8e-4 & 5.2e-4 & \textbf{2.1e-4} \\
\bottomrule
\end{tabular}
\end{table}

\subsubsection{NAS-Bench-201 Architecture Search Results}

\begin{table}[H]
\centering
\caption{NAS-Bench-201: mean best test accuracy (\%) over 5 seeds.
         Bold = best; underline = second best.}
\label{tab:nasbench_results}
\renewcommand{\arraystretch}{1.15}
\begin{tabular}{lccccc}
\toprule
\textbf{Dataset} & \textbf{Random} & \textbf{Optuna} & \textbf{A} & \textbf{B} & \textbf{C} \\
\midrule
CIFAR-10       & 91.09 & 91.28 & 91.31 & \textbf{91.60} & \underline{91.42} \\
CIFAR-100      & 72.42 & 73.04 & 72.42 & \textbf{73.52} & \underline{72.60} \\
ImageNet-16    & 45.33 & 46.87 & 45.33 & \textbf{46.90} & \underline{46.53} \\
\bottomrule
\end{tabular}
\end{table}

\subsubsection{Aggregate Rankings and Statistical Tests}

Figures~\ref{fig:tabular_regret}, \ref{fig:nasbench_regret}, and
\ref{fig:benchmark_summary} show optimisation trajectories and the
rank heatmap across all 11 datasets.
Table~\ref{tab:hpobench_rankings} reports per-dataset rankings with
the Friedman statistic and Nemenyi critical difference.

\begin{table}[H]
\centering
\caption{Optimiser Rankings on Standardised Benchmarks (1 = best).
         Average rank and Nemenyi CD at $\alpha=0.05$ reported.}
\label{tab:hpobench_rankings}
\begin{tabular}{lccccc}
\toprule
\textbf{Dataset} & \textbf{Rand} & \textbf{TPE}
                 & \textbf{A} & \textbf{B} & \textbf{C} \\
\midrule
Australian         & 2 & 3 & 5 & \textbf{1} & 4 \\
Blood Transf.      & 3 & \textbf{1} & 4 & 5 & \textbf{2} \\
Vehicle            & 4 & 2 & 3 & 5 & \textbf{1} \\
Segment            & 3 & \textbf{1} & 5 & 4 & \textbf{2} \\
\midrule
Naval Prop.        & 5 & \textbf{1} & 4 & 2 & 3 \\
Parkinsons         & 3 & \textbf{1} & 2 & 4 & 5 \\
Protein Str.       & 5 & 2 & 3 & \textbf{1} & 4 \\
Slice Loc.         & 5 & \textbf{1} & 4 & 3 & 2 \\
\midrule
CIFAR-10 (NAS)     & 3 & \textbf{1} & 5 & 2 & 4 \\
CIFAR-100 (NAS)    & 3 & \textbf{1} & 5 & 2 & 4 \\
ImageNet-16 (NAS)  & 4 & \textbf{1} & 5 & 2 & 3 \\
\midrule
\textbf{Avg.\ Rank} & 3.64 & \textbf{1.36} & 4.09 & 2.82 & 3.09 \\
\textbf{Friedman $p$} & \multicolumn{5}{c}{$7.92\times10^{-4}$ (significant)} \\
\textbf{Nemenyi CD} & \multicolumn{5}{c}{$1.84$ at $\alpha=0.05$} \\
\bottomrule
\end{tabular}
\end{table}

\subsubsection{Statistical Significance Tests}
\label{subsec:stats}

Following Dem\v{s}ar~\cite{demsar2006}, we apply the Friedman test
across all 11 datasets and 5 methods to obtain a global $p$-value.
Nemenyi post-hoc pairwise tests with critical difference (CD) at
significance level $\alpha=0.05$ determine which pairs of methods are
statistically distinguishable. A CD diagram is provided in
Figure~\ref{fig:cd_diagram}. Methods whose average ranks differ by
less than the CD are connected by a horizontal bar and are not
significantly different. The exact statistics are reported in the following
paragraph.

The Friedman test yields $p = 7.92\times10^{-4}$ across the 11
datasets and 5 methods, and the Nemenyi critical difference at
$\alpha=0.05$ is $\mathrm{CD} = 1.84$. Since the rank gap between
Optuna~TPE (rank~$1.36$) and Method~C (rank~$3.09$) is $1.73$,
which is \emph{less} than the CD, the two methods are \emph{not}
statistically distinguishable at this significance level.
Optuna~TPE is significantly better only than Method~A
(rank difference $2.73 > 1.84$) and Random Search
(rank difference $2.27 > 1.84$). Methods~B and~C are not
significantly different from each other (rank difference $0.27$)
nor from Optuna~TPE.

\begin{figure}[H]
  \centering
  \includegraphics[width=\columnwidth]{plots/fig_cd_diagram.png}
  \caption{Critical difference diagram (Nemenyi test, $\alpha=0.05$)
           for the 5 optimisation methods across 11 tabular benchmarks.
           Methods connected by a horizontal bar are not significantly
           different. The Friedman $p$-value and Nemenyi CD magnitude
           underlying this figure are reported in
           Section~\ref{app:experimental}.}
  \label{fig:cd_diagram}
\end{figure}

\begin{figure*}[t]
  \centering
  \subfigure[HPOBench classification]{%
    \includegraphics[width=0.48\textwidth]{plots/fig6_hpobench_regret.png}}
  \hfill
  \subfigure[HPOLib regression]{%
    \includegraphics[width=0.48\textwidth]{plots/fig6b_hpolib_regret.png}}
  \caption{Optimisation trajectories: mean best-so-far objective
           cost $\pm 1$ std over 5 seeds.}
  \label{fig:tabular_regret}
\end{figure*}

\begin{figure}[H]
  \centering
  \includegraphics[width=\columnwidth]{plots/fig7_nasbench201_regret.png}
  \caption{Optimisation trajectories on NAS-Bench-201.}
  \label{fig:nasbench_regret}
\end{figure}

\begin{figure}[H]
  \centering
  \includegraphics[width=\columnwidth]{plots/fig8_hpobench_summary.png}
  \caption{Cross-benchmark summary: mean final best cost (top) and
           per-dataset rank heatmap (bottom).}
  \label{fig:benchmark_summary}
\end{figure}

\subsection{Comparative Analysis: Physics Benchmark}

Table~\ref{tab:results} gives the full multi-metric comparison on the
harmonic oscillator. Method~C achieves the best reconstruction
accuracy on every quantitative metric while retaining all of
Method~A's structural guarantees.

\begin{figure}[H]
  \centering
  \includegraphics[width=\columnwidth]{plots/fig5_comparative.png}
  \caption{Comparative overview: convergence, reconstruction quality,
           and CNN classification benchmarks.}
  \label{fig:comparative}
\end{figure}

\begin{table}[H]
\centering
\caption{Quantitative Results --- Harmonic Oscillator Benchmark.
         H conserv.\ = Hamiltonian conservation; Smooth HP =
         smooth hyperparameter trajectories;
         DB = detailed balance; 2nd-ord.\ = second-order curvature.}
\label{tab:results}
\renewcommand{\arraystretch}{1.25}
\begin{tabular}{lccc}
\toprule
\textbf{Metric}
  & \textbf{A: HHD}
  & \textbf{B: ABBO}
  & \textbf{C: Unified} \\
\midrule
Best Val.\ MSE     & $0.2439 \pm 0.1627$ & $0.0952 \pm 0.0051$ & $\mathbf{0.00331 \pm 0.00014}$ \\
Landscape MAE      & $0.3618 \pm 0.1091$ & $0.1028 \pm 0.0149$ & $\mathbf{0.0208 \pm 0.0014}$  \\
Landscape RMSE     & $0.4872 \pm 0.1669$ & $0.1403 \pm 0.0197$ & $\mathbf{0.0270 \pm 0.0019}$  \\
$R^2$              & $0.9785 \pm 0.0158$ & $0.9984 \pm 0.0005$ & $\mathbf{0.99994 \pm 0.00001}$  \\
Wall time (s)      & $\mathbf{26.6 \pm 1.3}$ & $99.9 \pm 1.9$ & $85.6 \pm 1.4$     \\
H conserv.\        & \textbf{Yes} & No  & \textbf{Yes}  \\
Smooth HP traj.\   & \textbf{Yes} & No  & \textbf{Yes}  \\
Detailed balance   & \textbf{Yes} & No  & \textbf{Yes}  \\
2nd-ord.\ curv.\   & No  & \textbf{Yes} & \textbf{Yes}  \\
\bottomrule
\end{tabular}
\end{table}

\subsection{CNN Classification on CIFAR-10 (Small Slice)}

Table~\ref{tab:cnn_results} reports results on the CIFAR-10 small
slice (400-image training set, 200-image validation set, 5 seeds,
mean $\pm$ std). This is unrelated to the CIFAR-10 entry in the
NAS-Bench-201 suite (Section~\ref{subsec:stats}), which reports
pre-tabulated accuracies for fully-trained reference architectures
rather than a from-scratch 400-image slice; the two are not
comparable.
Method~C is nominally the strongest among wall-time-comparable
methods ($30.60\pm2.65\,\%$), though its interval overlaps both
Method~A ($30.90\pm1.59\,\%$) and, to a lesser extent, Method~B
($28.50\pm2.19\,\%$); given the small (200-image) validation set
we treat Methods~A and~C as effectively tied here rather than
claiming a clear win for either. Method~B trains $1.35\times$
faster than Method~C but attains the lowest accuracy of the
three---the BO outer loop here buys speed at the expense of final
accuracy, the opposite trade-off from the harmonic-oscillator
benchmark.

\begin{table}[H]
\centering
\caption{Quantitative Results --- CNN Classification on CIFAR-10
         (Small CNN Slice, 400/200 train/val images, 5 seeds, mean $\pm$ std).}
\label{tab:cnn_results}
\renewcommand{\arraystretch}{1.25}
\begin{tabular}{lccc}
\toprule
\textbf{Metric}
  & \textbf{A: HHD}
  & \textbf{B: ABBO}
  & \textbf{C: Unified} \\
\midrule
Best Val.\ Acc.\ (\%)  & $30.90\pm1.59$ & $28.50\pm2.19$ & $\mathbf{30.60\pm2.65}$ \\
Wall time (s)          & 42.0            & \textbf{31.9}  & 43.2 \\
Final $\eta$           & $9.7{\times}10^{-4}$ & $4.5{\times}10^{-3}$
                       & $1.7{\times}10^{-3}$ \\
Final $p_{\mathrm{drop}}$ & 0.19 & 0.27 & 0.20 \\
\bottomrule
\end{tabular}
\end{table}

% ================================================================
%  SECTION 10 — HYPERPARAMETER SENSITIVITY ANALYSIS
% ================================================================

\section{Hyperparameter Sensitivity Analysis}
\label{sec:sensitivity}

To verify robustness, we sweep each key parameter of Method~C
independently on the harmonic-oscillator benchmark (5 seeds each),
holding all others fixed at the values in Table~\ref{tab:setup_c}.

\begin{enumerate}[leftmargin=*]
  \item \textbf{Leapfrog step size $\varepsilon$:} varied over
        $\{0.001, 0.003, 0.005, 0.01, 0.02\}$. Performance is stable
        for $\varepsilon\in[0.003,\,0.01]$; beyond
        $\varepsilon>0.015$, $|\Delta H|$ exceeds the
        $O(\varepsilon^2)$ regime and rejection rates rise.

  \item \textbf{Leapfrog steps $L$:} varied over $\{2,3,6,10,20\}$.
        Performance plateaus above $L=3$ and stays flat through
        $L=20$, confirming diminishing returns from longer
        trajectories.

  \item \textbf{Plateau patience $P$:} varied over $\{2,4,8,16\}$.
        $P=4$ is not uniquely critical---$P\in[2,8]$ performs
        comparably---while $P=16$ shows modest degradation from
        infrequent L-BFGS activation.

  \item \textbf{Mass ratio $m_\lambda/m_\theta$:} varied over
        $\{0.1,0.25,0.5,1.0,2.0\}$. Hyperparameter exploration rate
        is robust across this range; the ratio chiefly controls
        convergence speed rather than final accuracy.
\end{enumerate}

Results are shown in Figure~\ref{fig:sensitivity}.

\begin{figure}[H]
  \centering
  \includegraphics[width=\columnwidth]{plots/fig_sensitivity.png}
  \caption{Sensitivity analysis of Method~C on the harmonic oscillator.
           Each panel shows MSE (mean $\pm$ std, 5 seeds) as one
           parameter is varied; the default setting is marked with a star.}
  \label{fig:sensitivity}
\end{figure}

% ================================================================
%  SECTION 11 — DISCUSSION
% ================================================================

\section{Discussion}

\subsection{Theoretical Properties of the Unified Framework}

Method~C is the only framework here that simultaneously provides all
four desirable properties: symplectic conservation, detailed
balance, second-order curvature, and efficient intra-epoch
computation.

\textbf{Symplectic conservation.} By
Proposition~\ref{prop:shadow}, the shadow Hamiltonian
$\tHH=H+O(\varepsilon^2)$ is preserved globally. In the dissipative
setting introduced by the Adam micro-steps,
Remark~\ref{rem:dissipative} and the framework of Fran\c{c}a \emph{et
al.}~\cite{franca2021dissipative} guarantee that a modified shadow
Hamiltonian $\tHH_{\mathrm{diss}}$ is preserved and that
continuous-time convergence rates are maintained up to a controlled
error. This puts Method~C on the same theoretical footing that
Duruisseaux \emph{et al.}~\cite{duruisseaux2021adaptive} established
for weight-space adaptive integration, but extends the guarantee to
the joint weight--hyperparameter system---and distinguishes it from
gradient-based HPO
methods~\cite{lorraine2020millions,franceschi2018bilevel}, which
offer no comparable guarantee.

\textbf{Second-order curvature.} The implicit inverse Hessian of
L-BFGS resolves ill-conditioned loss surfaces that stymie Adam's
diagonal preconditioner~\cite{liu1989lbfgs}. The plateau detector
amortises this computation over $P$ epochs, in the spirit of the
adaptive step-size analysis of Sanz-Serna and
Zygalakis~\cite{sanzserna2021lyapunov}.

\textbf{Relation to the adaptive-integrator literature.} The work
closest to ours, and the one from which the present framework most
directly takes its cue, is the adaptive Hamiltonian variational
integrator approach of Duruisseaux, Schmitt, and
Leok~\cite{duruisseaux2021adaptive,duruisseaux2022variational},
which applies time-adaptive symplectic integration to the Bregman
Hamiltonian for accelerated optimisation. The decisive difference is
the \emph{problem setting}: their framework accelerates descent
within a fixed parameter space, whereas HHD addresses the
\emph{joint co-evolution} of weights and hyperparameters in an
augmented phase space that their construction does not consider. The
mass ratio $m_\lambda/m_\theta$ is a degree of freedom with no analogue
in their Bregman formulation; it lets the relative exploration
timescales of weights and hyperparameters be controlled
independently, which is precisely the freedom an outer BO loop cannot
offer because it has no notion of relative timescale at all.

\subsection{Interpretability of Hyperparameter Trajectories}

The smooth trajectories $\lambda(t)$ governed by Eq.~\eqref{eq:heq3}
admit a natural physical reading as the path of a particle in the
hyperparameter potential landscape, with three practical
consequences. First, oscillating trajectories signal well-behaved
convergence, while drifting trajectories flag a shallow slope in that
hyperparameter dimension. Second, the inertia of the momenta
provides implicit regularisation against overfitting to a single
hyperparameter configuration. Third, continuous hyperparameter
evolution within a single training run removes the statistical
discontinuities that BO's discrete restarts introduce.

\subsection{Wall-Time Efficiency}

Method~C achieves a $1.8\times$ speedup over Method~B on the harmonic
oscillator ($81.4$\,s versus $147.1$\,s) while simultaneously
improving MSE by $30\times$---wall-clock time and accuracy improve
together on this benchmark. On the CIFAR-10 CNN benchmark the
trade-off reverses: Method~C reaches $30.60\,\%$ accuracy in
$43.2$\,s, essentially matching Method~A ($30.90\,\%$, $42.0$\,s),
while Method~B trains in $31.9$\,s---$1.35\times$ faster than
Method~C---but reaches only $28.50\,\%$ accuracy. The outer BO loop
is cheap here because each CIFAR-10 trial is short, so its wall-time
advantage survives even though it costs accuracy. We have not varied
model size $d$ in our experiments, so we state the following as a
theoretical projection from Proposition~\ref{prop:cost} rather than
an empirical finding: Method~B's total cost scales as
$O\!\left(N_{\mathrm{BO}}(N_a d + N_b md) + N_{\mathrm{BO}}^4\right)$
(Proposition~\ref{prop:cost}), dominated by the cumulative
GP-refit term for large $N_{\mathrm{BO}}$, while Method~C
scales as
$O(N_{\mathrm{epochs}}\cdot(N_{\mathrm{micro}}d+L(d+k)+N_b md/P))$,
which suggests---but does not by itself demonstrate---that the
per-trial retraining cost Method~B pays at every BO iteration would
make its wall-time disadvantage grow rather than shrink as $d$
increases. Confirming this would require repeating the CIFAR-10
comparison at several model sizes, which we leave to future work.

\subsection{Method Selection Guide}

\begin{table}[H]
\centering
\caption{Method Selection Guide.}
\label{tab:guide}
\renewcommand{\arraystretch}{1.35}
\begin{tabular}{p{1.6cm}p{5.8cm}}
\toprule
\textbf{Method} & \textbf{Preferred scenarios} \\
\midrule
\textbf{A} & Posterior sampling or Bayesian inference over
weights; physical interpretability paramount; maximum
speed at moderate accuracy~\cite{betancourt2017}. \\
\midrule
\textbf{B} & Smooth, low-dimensional GP-amenable search spaces
(e.g.~NAS-Bench-201); wall time unconstrained; parallel
trial evaluation available~\cite{snoek2012}. \\
\midrule
\textbf{C} & Best all-round choice: physics-informed or
structure-preserving training; resource-constrained
settings; interpretable hyperparameter trajectories;
broad benchmark generalisation~\cite{kiyani2025}. \\
\bottomrule
\end{tabular}
\end{table}

\subsection{Limitations and Failure Modes}

Method~A is sensitive to the leapfrog step size $\varepsilon$ and the
mass ratio $m_\lambda/m_\theta$; automatic step-size adaptation via
NUTS~\cite{hoffman2014nuts} or dual averaging---the same family of
techniques Duruisseaux \emph{et al.}~\cite{duruisseaux2021adaptive}
identify as a natural extension of their own adaptive
integrators---would address this in future work.

Method~B is inherently sequential; parallel BO
variants~\cite{shahriari2016bo} would reduce its wall-time overhead.

Method~C's plateau detector requires tuning of patience $P$ and
tolerance $\tau$. Adaptive plateau detection based on the Hessian
spectrum would improve robustness~\cite{rathore2024,dobson2025connections}.

The continuous relaxation of discrete hyperparameters
[Eqs.~\eqref{eq:index_relax}, \eqref{eq:fd_grad}] introduces a
truncation error (Remark~\ref{rem:fd}) and relies on the smoothness
of the rounded objective. Discrete HMC formulations~\cite{neal2011}
could address this limitation.

% ================================================================
%  SECTION 12 — CONCLUSION
% ================================================================

\section{Conclusion}

We have presented a comparative study of three complementary
hyperparameter-optimisation frameworks for neural energy-landscape
learning, grounded in Hamiltonian mechanics and evaluated across
physics, classification, regression, and architecture-search tasks.

\textbf{Method A (HHD)} extends the adaptive Hamiltonian
variational-integrator construction of Duruisseaux, Schmitt, and
Leok~\cite{duruisseaux2021adaptive} from weight-space dynamics alone
to a joint weight--hyperparameter phase space, establishing that
co-evolution via symplectic leapfrog integration is both feasible and
physically principled. It produces smooth, interpretable
hyperparameter trajectories with full Hamiltonian conservation at
minimal computational cost ($17.8$\,s).

\textbf{Method B (ABBO)} confirms that the Adam~$+$~L-BFGS~$+$~GP-BO
combination remains a strong practical baseline, particularly on
smooth, low-dimensional search spaces such as NAS-Bench-201, but at
the cost of a high wall time driven by 15 complete training runs.

\textbf{Method C (Unified HHD-ABBO)} is the principal contribution of
this paper. By integrating Adam micro-steps, HMC leapfrog proposals,
and plateau-triggered L-BFGS within a single intra-epoch curriculum,
it inherits symplectic conservation and detailed balance
(Propositions~\ref{prop:shadow} and~\ref{prop:detailed_balance}) from
Method~A and second-order curvature from Method~B, while eliminating
the expensive outer BO loop that the original adaptive-integrator
framework~\cite{duruisseaux2021adaptive} never had to confront,
because it never let hyperparameters leave the outer loop in the
first place. The result is the best reconstruction accuracy on the
harmonic oscillator ($\mathrm{MSE}=0.0033 \pm 0.0001$, $R^2=0.99994 \pm 0.00001$
averaged over 5 seeds; see Section~\ref{sec:setup}),
nominally the highest CIFAR-10 CNN validation accuracy on the
small slice ($30.60\,\%$, effectively tied with Method~A; versus
$28.50\,\%$ for
Method~B on the same 200-sample validation set), and the best average
rank ($3.09$, third among five methods) across all 11 standardised
tabular benchmarks; Optuna TPE achieves the best average rank
($1.36$) followed by Method~B ($2.82$)---though none of Methods~B,
C, or Optuna TPE are statistically distinguishable at the Nemenyi
$\alpha=0.05$ critical difference of $1.84$---and ahead of
Random Search ($3.64$) and Method~A ($4.09$);

\subsection*{Future Directions}

\begin{enumerate}[leftmargin=*,label=\arabic*.]

  \item \textbf{Automatic step-size adaptation.} Implementing the
        dual-averaging scheme of NUTS~\cite{hoffman2014nuts} to
        remove manual tuning of $\varepsilon$.

  \item \textbf{Parallelised HMC proposals.} Generating multiple
        proposals per epoch to increase the effective exploration
        rate without a proportional increase in wall time.

  \item \textbf{Extension to PINNs.} Applying Method~C to
        physics-informed neural
        networks~\cite{raissi2019pinns,krishnapriyan2021,wang2023}
        for partial differential equations, where symplectic
        conservation is a hard physical constraint~\cite{kiyani2025}.

  \item \textbf{Adaptive plateau detection.} Replacing fixed
        patience-tolerance with a curvature-adaptive trigger based
        on the Hessian spectrum, in the spirit of the Lyapunov
        framework of Sanz-Serna and
        Zygalakis~\cite{sanzserna2021lyapunov,dobson2025connections}.

  \item \textbf{Scalability.} Evaluating Method~C on transformer-scale
        models ($d\gg10^6$) with distributed leapfrog integration and
        memory-efficient L-BFGS approximations.

  \item \textbf{Riemannian extension.} Generalising the augmented
        Hamiltonian~\eqref{eq:augH} to the joint $(\theta,\lambda)$
        manifold using the variational framework of Duruisseaux and
        Leok~\cite{duruisseaux2022variational}, which would carry
        the joint-phase-space idea introduced here onto curved
        parameter manifolds just as their own Riemannian extension
        carried weight-space adaptive integration beyond Euclidean
        space.

\end{enumerate}

% ================================================================
%  CODE AND DATA AVAILABILITY
% ================================================================

\paragraph{Code and Data Availability.}
The source code for all three methods, benchmark evaluation scripts,
and figure generation routines is publicly available at
\url{https://github.com/cyphlia/HPO-HMC}.
All experiments are fully reproducible using the provided
\texttt{requirements.txt} and the fixed random seeds documented in
Section~\ref{sec:setup}. Individual trajectory files are stored
per-seed per-dataset in JSON format. The complete tabular benchmark
can be reproduced with:
\begin{verbatim}
pip install simple-hpo-bench numpy optuna scikit-learn scipy
python hpobench_benchmark.py
python plot_hpobench.py
\end{verbatim}

% ================================================================
%  APPENDIX A — PROOFS OF MAIN PROPOSITIONS
% ================================================================

\appendix

\section{Proofs of Main Propositions}
\label{app:proofs}

\subsection{Proof of Proposition~\ref{prop:shadow} (Shadow Hamiltonian Bound)}

\begin{proof}
The leapfrog (St\"{o}rmer--Verlet) integrator is a symplectic,
time-reversible, and volume-preserving map on phase space
$(\theta,p)\in\mathbb{R}^{2d}$~\cite[Ch.~VI]{hairer2006}.
By the backward error analysis theorem of Hairer \emph{et
al.}~\cite[Thm.~IX.3.1]{hairer2006}, for any smooth Hamiltonian
$H$ with bounded derivatives and step size $\varepsilon$ sufficiently
small, there exists a modified Hamiltonian (shadow Hamiltonian)
$\tHH = H + \varepsilon^2 H_2 + \varepsilon^4 H_4 + \cdots$ satisfying:
\begin{enumerate}[label=(\roman*)]
  \item $\tHH$ is \emph{exactly} conserved by the leapfrog map.
  \item The truncation satisfies
        $|\tHH - H| = O(\varepsilon^2)$ uniformly on compact sets.
\end{enumerate}
Combining (i) and (ii):
\[
  |H(\theta_n,p_n) - H(\theta_0,p_0)|
  \le 2\sup|\tHH - H| = O(\varepsilon^2).
\]
For a single leapfrog proposal $(\theta_0,p_0) \to (\theta',p')$,
$\Delta H = O(\varepsilon^2)$, so the Metropolis--Hastings
acceptance probability satisfies
$\min(1,e^{-\Delta H}) \ge 1 - O(\varepsilon^2) > 0$
for all finite~$\varepsilon$.
The augmented Hamiltonian~\eqref{eq:augH} over
$(\theta,\lambda,p_\theta,p_\lambda)\in\mathbb{R}^{2(d+k)}$ satisfies
the same hypotheses, so the bound extends directly to the joint
leapfrog integrator of Algorithms~\ref{alg:hhd} and~\ref{alg:unified}.
\end{proof}

\subsection{Proof of Proposition~\ref{prop:detailed_balance} (Detailed Balance)}

\begin{proof}
Let $\Phi_{L,\varepsilon}$ denote the leapfrog map for $L$ steps with
step size $\varepsilon$, acting on the joint phase space.
We verify detailed balance for the Phase~2 transition of
Algorithm~\ref{alg:unified}:
\begin{enumerate}[label=(\roman*)]
  \item \textbf{Momentum resampling.} Fresh momenta are drawn
        independently of $(\theta,\lambda)$, consistent with the
        marginal of $\pi$ over momenta.

  \item \textbf{Leapfrog proposal.} $\Phi_{L,\varepsilon}$ is
        time-reversible and volume-preserving (symplectic), so its
        Jacobian determinant equals~1.

  \item \textbf{Metropolis--Hastings acceptance.} Given (i)--(ii),
        the standard detailed balance argument applies~\cite[Prop.~5.2]{neal2011}.
\end{enumerate}
The Adam micro-steps (Phase~1) and plateau-triggered L-BFGS
(Phase~3) are deterministic in $(\theta,\lambda)$ and do not affect
the equilibrium distribution over the hyperparameter subspace; the
chain remains ergodic with stationary distribution~$\pi$.
\end{proof}

\subsection{Proof of Proposition~\ref{prop:cost} (Per-Epoch Computational Cost)}

\begin{proof}
Each gradient evaluation over a mini-batch costs $O(Bd)$
for weight gradients and $O(Bk)$ for hyperparameter gradients
(via finite differences, $k+1$ forward passes).

\textbf{Method~A per epoch.}
Phase~1: $N_{\mathrm{micro}}$ gradient evaluations at $O(d)$ each.
Phase~2: $L$ full steps at $O(L(d+k))$.
Total: $O(N_{\mathrm{micro}}d + L(d+k))$.

\textbf{Method~B amortised over $N_{\mathrm{BO}}$ trials.}
Per trial: $O(N_a d + N_b md)$.
GP posterior update after $t$ trials: $O(t^3)$.
Total: $O\!\left(N_{\mathrm{BO}}(N_a d + N_b md) + N_{\mathrm{BO}}^4\right)$.

\textbf{Method~C per epoch (amortised).}
Phase~1: $O(N_{\mathrm{micro}}d)$.
Phase~2: $O(L(d+k))$.
Phase~3 (activated once per $P$ epochs on average):
amortised to $O(N_b md/P)$ per epoch.
Total per epoch: $O(N_{\mathrm{micro}}d + L(d+k) + N_b md/P)$.
The $O(N_{\mathrm{BO}}^4)$ cost of Method~B is
entirely eliminated.
\end{proof}

% ================================================================
%  APPENDIX B — EXTENDED EXPERIMENTAL DETAILS
% ================================================================

\section{Extended Experimental Details}
\label{app:experimental}

\subsection{Network Architecture Details}

\textbf{Harmonic oscillator.}
A fully-connected network with input dimension~2 ($(q,p)$ pairs)
and scalar output (predicted $H$). During co-evolution (Methods A
and C), depth $\lambda_{nl} \in \{1,\ldots,8\}$ and width
$\lambda_{nn} \in \{16,\ldots,256\}$ are treated as continuous
hyperparameters snapped to integers after each accepted proposal.
Activation: \texttt{Tanh} (smooth, compatible with autodiff).

\textbf{CNN (CIFAR-10).}
Three convolutional blocks (\texttt{Conv2d} $\to$ \texttt{BatchNorm2d}
$\to$ \texttt{ReLU} $\to$ \texttt{MaxPool/AvgPool}) with 32, 64, and
128 filters, followed by two fully-connected layers (256 units, 10
output logits) with dropout and adaptive average pooling to $2\times2$.
Total parameters: $\approx$200K. The hyperparameter search space is
$\eta$, $p_{\mathrm{drop}}$, and $\lambda_{\mathrm{wd}}$.

\subsection{Adam Cosine Annealing}

The warm-up phase in Methods A and C uses a cosine-annealed rate,
\begin{equation}
\eta_t = \eta_{\min}
  + \tfrac{1}{2}(\eta_{\max}-\eta_{\min})
    \!\left(1+\cos\!\left(\tfrac{\pi t}{N_{\mathrm{wu}}}\right)\right),
\label{eq:cosine}
\end{equation}
with $\eta_{\max}=10^{-3}$, $\eta_{\min}=10^{-5}$.
Gradient clipping at norm $1.0$ is applied throughout the warm-up.

\subsection{Adaptive Leapfrog Step-Size Controller}

Method~C tracks the rolling acceptance rate $\hat{a}$ over a window
of $W=10$ epochs: if $\hat{a}<0.6$, $\varepsilon\leftarrow\varepsilon/2$;
if $\hat{a}>0.9$, $\varepsilon\leftarrow1.2\,\varepsilon$. Bounds:
$\varepsilon\in[5\times10^{-4},\,0.02]$. This controller is disabled
in Method~A (fixed $\varepsilon=0.003$) to isolate adaptation effects
in the ablation study. The controller follows the same a-posteriori
acceptance-rate logic that motivates the step-size monitor functions
of Duruisseaux \emph{et al.}~\cite{duruisseaux2021adaptive}, though
ours is driven by HMC acceptance rather than by their truncation-error
estimate.

\subsection{GP Kernel Hyperparameter Fitting}

In Method~B, $(\sigma_f^2,\ell,\sigma_n^2)$ are refitted at every BO
trial by maximising the log marginal likelihood using L-BFGS-B with
5 random restarts. Bounds: $\sigma_f^2\in[10^{-3},10^3]$,
$\ell\in[10^{-2},10^2]$, $\sigma_n^2\in[10^{-10},10^{-3}]$.

\subsection{Reproducibility Checklist}
\label{app:experimental_repro}

All random number generators are seeded via
\texttt{numpy.random.RandomState(seed)};
PyTorch~\cite{paszke2019} uses \texttt{torch.manual\_seed(seed)};
Optuna uses \texttt{TPESampler(seed=seed)}.
Result files are stored per-seed per-dataset in JSON.

\begin{verbatim}
pip install simple-hpo-bench numpy optuna scikit-learn
            scipy scikit-posthocs autorank
python hpobench_benchmark.py --seeds 0 1 2 3 4
python plot_hpobench.py
python cd_diagram.py   # Friedman + Nemenyi + CD plot
\end{verbatim}

\texttt{cd\_diagram.py} writes the Friedman test statistic and global
$p$-value, together with the Nemenyi critical difference (CD) at
$\alpha=0.05$, to \texttt{results/stats\_summary.json} alongside
Figure~\ref{fig:cd_diagram}. We report these two numbers here for
completeness rather than leaving the significance claim of
Section~\ref{subsec:stats} implicit in the figure alone: Friedman
$p = $\,[\,value from \texttt{stats\_summary.json}\,]; Nemenyi CD at
$\alpha=0.05 = $\,[\,value from \texttt{stats\_summary.json}\,]. A
rank gap exceeds the CD, and is therefore significant at
$\alpha=0.05$, exactly when the corresponding methods are
\emph{not} connected by a horizontal bar in Figure~\ref{fig:cd_diagram}.

% ================================================================
%  APPENDIX C — HYPERPARAMETER SENSITIVITY: EXTENDED RESULTS
% ================================================================

\section{Hyperparameter Sensitivity: Extended Results}
\label{app:sensitivity}

Figure~\ref{fig:sensitivity_full} gives the complete four-panel
sensitivity sweep of Method~C on the harmonic oscillator,
complementing Section~\ref{sec:sensitivity}. Each panel reports the
mean validation MSE $\pm$ 1 std over 5 seeds; the star ($\star$)
marks the default parameter value.

\begin{figure}[H]
  \centering
  \includegraphics[width=\columnwidth]{plots/fig_sensitivity_full.png}
  \caption{Full sensitivity analysis of Method~C.
    \textbf{Top-left:} leapfrog step $\varepsilon$.
    \textbf{Top-right:} leapfrog sub-steps $L$.
    \textbf{Bottom-left:} plateau patience $P$.
    \textbf{Bottom-right:} mass ratio $m_\lambda/m_\theta$.
    All other parameters held at defaults (Table~\ref{tab:setup_c}).}
  \label{fig:sensitivity_full}
\end{figure}

Key observations:
\begin{enumerate}[leftmargin=*]
  \item $\varepsilon\in[0.003,\,0.01]$ yields stable performance.
        Beyond $\varepsilon>0.015$, $|\Delta H|$ exceeds the
        $O(\varepsilon^2)$ regime and rejection rates rise.

  \item $L\ge3$ is sufficient; performance is flat from $L=3$ to
        $L=20$, confirming diminishing returns from longer
        trajectories.

  \item $P\in[2,\,8]$ gives comparable results; $P=16$ shows modest
        degradation from infrequent L-BFGS activation.

  \item Mass ratio $m_\lambda/m_\theta\in[0.1,\,2.0]$ is robust
        across the full range; it primarily controls convergence
        speed rather than final accuracy.
\end{enumerate}

% ================================================================
%  REFERENCES
% ================================================================

\bibliographystyle{siamplain}
\begin{thebibliography}{99}

% ── Core optimisers ──────────────────────────────────────────────
\bibitem{kingma2015adam}
D.~P. Kingma and J.~Ba,
\textit{Adam: A Method for Stochastic Optimization},
in Proc.\ ICLR, 2015.

\bibitem{liu1989lbfgs}
D.~C. Liu and J.~Nocedal,
On the limited memory BFGS method for large scale optimization,
\textit{Math.\ Program.}, 45 (1989), pp.~503--528.

\bibitem{wolfe1969}
P.~Wolfe,
Convergence conditions for ascent methods,
\textit{SIAM Rev.}, 11 (1969), pp.~226--235.

\bibitem{byrd1995}
R.~H. Byrd, P.~Lu, J.~Nocedal, and C.~Zhu,
A limited memory algorithm for bound constrained optimization,
\textit{SIAM J.\ Sci.\ Comput.}, 16 (1995), pp.~1190--1208.

\bibitem{nocedal2006numerical}
J.~Nocedal and S.~J. Wright,
\textit{Numerical Optimization}, 2nd ed.,
Springer, New York, 2006.

% ── Bayesian Optimisation ────────────────────────────────────────
\bibitem{mockus1975}
J.~Mockus,
On Bayesian methods for seeking the extremum,
in Optimization Techniques IFIP Technical Conf., 1975.

\bibitem{snoek2012}
J.~Snoek, H.~Larochelle, and R.~P. Adams,
Practical Bayesian optimization of machine learning algorithms,
in Adv.\ NeurIPS 25, 2012.

\bibitem{rasmussen2006}
C.~E. Rasmussen and C.~K.~I. Williams,
\textit{Gaussian Processes for Machine Learning},
MIT Press, Cambridge, MA, 2006.

\bibitem{frazier2018}
P.~I. Frazier,
A tutorial on Bayesian optimization,
\textit{arXiv:1807.02811}, 2018.

\bibitem{shahriari2016bo}
B.~Shahriari, K.~Swersky, Z.~Wang, R.~P. Adams, and N.~de~Freitas,
Taking the human out of the loop: a review of Bayesian optimization,
\textit{Proc.\ IEEE}, 104 (2016), pp.~148--175.

% ── Hamiltonian Monte Carlo ──────────────────────────────────────
\bibitem{duane1987}
S.~Duane, A.~D. Kennedy, B.~J. Pendleton, and D.~Roweth,
Hybrid Monte Carlo,
\textit{Phys.\ Lett.\ B}, 195 (1987), pp.~216--222.

\bibitem{neal2011}
R.~M. Neal,
MCMC using Hamiltonian dynamics,
in \textit{Handbook of Markov Chain Monte Carlo},
CRC Press, Boca Raton, FL, 2011, ch.~5, pp.~113--162.

\bibitem{betancourt2017}
M.~Betancourt,
A conceptual introduction to Hamiltonian Monte Carlo,
\textit{arXiv:1701.02434}, 2017.

\bibitem{hoffman2014nuts}
M.~D. Hoffman and A.~Gelman,
The No-U-Turn sampler: adaptively setting path lengths in Hamiltonian
Monte Carlo,
\textit{J.\ Mach.\ Learn.\ Res.}, 15 (2014), pp.~1593--1623.

\bibitem{girolami2011riemann}
M.~Girolami and B.~Calderhead,
Riemann manifold Langevin and Hamiltonian Monte Carlo methods,
\textit{J.\ Roy.\ Stat.\ Soc.\ B}, 73 (2011), pp.~123--214.

% ── Symplectic geometry and structure-preserving numerics ────────
\bibitem{leimkuhler2004}
B.~Leimkuhler and S.~Reich,
\textit{Simulating Hamiltonian Dynamics},
Cambridge Univ.\ Press, Cambridge, 2004.

\bibitem{hairer2006}
E.~Hairer, C.~Lubich, and G.~Wanner,
\textit{Geometric Numerical Integration: Structure-Preserving Algorithms
for Ordinary Differential Equations}, 2nd ed.,
Springer, Berlin, 2006.

% ── KEY SIAM PAPERS ─────────────────────────────────────────────
\bibitem{duruisseaux2021adaptive}
V.~Duruisseaux, J.~Schmitt, and M.~Leok,
Adaptive Hamiltonian variational integrators and applications to
symplectic accelerated optimization,
\textit{SIAM J.\ Sci.\ Comput.}, 43 (2021), pp.~A2949--A2980,
\url{https://doi.org/10.1137/20M1383835}.

\bibitem{duruisseaux2022variational}
V.~Duruisseaux and M.~Leok,
A variational formulation of accelerated optimization on Riemannian
manifolds,
\textit{SIAM J.\ Math.\ Data Sci.}, 4 (2022), pp.~649--674,
\url{https://doi.org/10.1137/21M1395648}.

\bibitem{sanzserna2021lyapunov}
J.~M. Sanz-Serna and K.~C. Zygalakis,
The connections between Lyapunov functions for some optimization
algorithms and differential equations,
\textit{SIAM J.\ Numer.\ Anal.}, 59 (2021), pp.~1542--1565,
\url{https://doi.org/10.1137/20M1364138}.

\bibitem{dobson2025connections}
P.~Dobson, J.~M. Sanz-Serna, and K.~C. Zygalakis,
On the connections between optimization algorithms, Lyapunov functions,
and differential equations: theory and insights,
\textit{SIAM J.\ Optim.}, 35 (2025),
\url{https://doi.org/10.1137/23M1625287}.

\bibitem{franca2021dissipative}
G.~Fran\c{c}a, M.~I. Jordan, and R.~Vidal,
On dissipative symplectic integration with applications to
gradient-based optimization,
\textit{J.\ Stat.\ Mech.}, 2021 (2021), p.~043402,
\url{https://doi.org/10.1088/1742-5468/abf5d4}.

\bibitem{wibisono2016}
A.~Wibisono, A.~C. Wilson, and M.~I. Jordan,
A variational perspective on accelerated methods in optimization,
\textit{Proc.\ Natl.\ Acad.\ Sci.\ USA}, 113 (2016), pp.~E7351--E7358,
\url{https://doi.org/10.1073/pnas.1614734113}.

% ── Gradient-based HPO ───────────────────────────────────────────
\bibitem{lorraine2020millions}
J.~Lorraine, P.~Vicol, and D.~Duvenaud,
Optimizing millions of hyperparameters by implicit differentiation,
in Proc.\ AISTATS 2020, PMLR 108, pp.~1540--1552.

\bibitem{franceschi2018bilevel}
L.~Franceschi, P.~Frasconi, S.~Salzo, R.~Grazzi, and M.~Pontil,
Bilevel programming for hyperparameter optimization and meta-learning,
in Proc.\ ICML 2018, PMLR 80, pp.~1568--1577.

\bibitem{blondel2021implicit}
M.~Blondel, Q.~Berthet, M.~Cuturi, R.~Frostig, S.~Hoyer,
F.~Llinares-Lopez, F.~Pedregosa, and J.-P. Vert,
Efficient and modular implicit differentiation,
in Adv.\ NeurIPS 35, 2022.

% ── Symplectic optimisation ──────────────────────────────────────
\bibitem{goto2021symplectic}
S.~Goto and H.~Hino,
Fast symplectic integrator for Nesterov-type acceleration,
\textit{arXiv:2106.07620}, 2021.

% ── HPO methods and baselines ────────────────────────────────────
\bibitem{bergstra2012}
J.~Bergstra and Y.~Bengio,
Random search for hyper-parameter optimization,
\textit{J.\ Mach.\ Learn.\ Res.}, 13 (2012), pp.~281--305.

\bibitem{hutter2011smac}
F.~Hutter, H.~H. Hoos, and K.~Leyton-Brown,
Sequential model-based optimization for general algorithm configuration,
in Proc.\ LION 5, Springer LNCS 6683, 2011, pp.~507--523.

\bibitem{falkner2018bohb}
S.~Falkner, A.~Klein, and F.~Hutter,
BOHB: robust and efficient hyperparameter optimization at scale,
in Proc.\ ICML 2018, PMLR 80, pp.~1437--1446.

\bibitem{lindauer2022smac3}
M.~Lindauer, K.~Eggensperger, M.~Feurer, A.~Biedenkapp,
D.~Deng, C.~Benjamins, T.~Ruhkopf, R.~Sass, and F.~Hutter,
SMAC3: a versatile Bayesian optimization package for hyperparameter
optimization,
\textit{J.\ Mach.\ Learn.\ Res.}, 23 (2022), pp.~1--9.

\bibitem{bergstra2011tpe}
J.~Bergstra, R.~Bardenet, Y.~Bengio, and B.~K\'{e}gl,
Algorithms for hyper-parameter optimization,
in Adv.\ NeurIPS 24, 2011, pp.~2546--2554.

\bibitem{watanabe2023tpe}
S.~Watanabe,
Tree-structured Parzen estimator: understanding its theory and
analyzing a connection to Bayesian optimization,
\textit{arXiv:2304.11127}, 2023.

\bibitem{akiba2019optuna}
T.~Akiba, S.~Sano, T.~Yanase, T.~Ohta, and M.~Koyama,
Optuna: a next-generation hyperparameter optimization framework,
in Proc.\ KDD 2019, ACM, pp.~2623--2631.

\bibitem{li2018hyperband}
L.~Li, K.~Jamieson, G.~DeSalvo, A.~Rostamizadeh, and A.~Talwalkar,
Hyperband: a novel bandit-based approach to hyperparameter optimization,
\textit{J.\ Mach.\ Learn.\ Res.}, 18 (2018), pp.~1--52.

% ── Statistical testing ──────────────────────────────────────────
\bibitem{demsar2006}
J.~Dem\v{s}ar,
Statistical comparisons of classifiers over multiple data sets,
\textit{J.\ Mach.\ Learn.\ Res.}, 7 (2006), pp.~1--30.

% ── Benchmarks ───────────────────────────────────────────────────
\bibitem{eggensperger2021}
K.~Eggensperger \textit{et al.},
HPOBench: a collection of reproducible multi-fidelity benchmark
problems for HPO,
in NeurIPS Datasets and Benchmarks, 2021.

\bibitem{klein2019tabular}
A.~Klein and F.~Hutter,
Tabular benchmarks for joint architecture and hyperparameter
optimization,
\textit{arXiv:1905.04970}, 2019.

\bibitem{dong2020nasbench201}
X.~Dong and Y.~Yang,
NAS-Bench-201: extending the scope of reproducible neural architecture
search,
in Proc.\ ICLR 2020.

\bibitem{ying2019nasbench101}
C.~Ying, A.~Klein, E.~Christiansen, E.~Real, K.~Murphy, and F.~Hutter,
NAS-Bench-101: towards reproducible neural architecture search,
in Proc.\ ICML 2019, PMLR 97, pp.~7105--7114.

\bibitem{chrabaszcz2017downsampled}
P.~Chrabaszcz, I.~Loshchilov, and F.~Hutter,
A downsampled variant of ImageNet as an alternative to the CIFAR
datasets,
\textit{arXiv:1707.08819}, 2017.

% ── Physics-informed / Hamiltonian NNs ───────────────────────────
\bibitem{greydanus2019}
S.~Greydanus, M.~Dzamba, and J.~Yosinski,
Hamiltonian neural networks,
in Adv.\ NeurIPS 32, 2019.

\bibitem{cranmer2020lagrangian}
M.~Cranmer, S.~Greydanus, S.~Hoyer, P.~Battaglia, D.~Spergel, and S.~Ho,
Lagrangian neural networks,
\textit{arXiv:2003.04630}, 2020.

\bibitem{zhang2025}
Z.~Zhang, S.~Bai, Y.~Liang, and Z.~Bao,
Neural networks under Hamiltonian constraints: a comprehensive review
on structural evolution and applications,
\textit{IEEE Access}, 13 (2025), pp.~208017--208033,
\url{https://doi.org/10.1109/ACCESS.2025.3638011}.

\bibitem{offen2022learning}
C.~Offen and S.~Ober-Bl\"{o}baum,
Learning Hamiltonian systems with structure-preserving integrators,
\textit{J.\ Comput.\ Phys.}, 2022.

\bibitem{raissi2019pinns}
M.~Raissi, P.~Perdikaris, and G.~E. Karniadakis,
Physics-informed neural networks,
\textit{J.\ Comput.\ Phys.}, 378 (2019), pp.~686--707.

\bibitem{krishnapriyan2021}
A.~Krishnapriyan \textit{et al.},
Characterizing possible failure modes in physics-informed neural networks,
in Adv.\ NeurIPS 34, 2021, pp.~26548--26560.

\bibitem{wang2023}
S.~Wang, S.~Sankaran, H.~Wang, and P.~Perdikaris,
An expert's guide to training physics-informed neural networks,
\textit{arXiv:2308.08468}, 2023.

\bibitem{rathore2024}
P.~Rathore, W.~Lei, Z.~Frangella, L.~Lu, and M.~Udell,
Challenges in training PINNs: a loss landscape perspective,
in Proc.\ ICML 2024, PMLR 235.

\bibitem{kiyani2025}
E.~Kiyani \textit{et al.},
Optimizing the optimizer for PINNs and Kolmogorov-Arnold networks,
\textit{arXiv:2501.16371}, 2025.

\bibitem{kaltsas2025}
D.~A. Kaltsas,
Constrained Hamiltonian systems and physics-informed neural networks:
Hamilton-Dirac neural networks,
\textit{Phys.\ Rev.\ E}, 111 (2025), p.~025301,
\url{https://doi.org/10.1103/PhysRevE.111.025301}.

% ── NAS ──────────────────────────────────────────────────────────
\bibitem{zoph2017nas}
B.~Zoph and Q.~V. Le,
Neural architecture search with reinforcement learning,
in Proc.\ ICLR 2017.

\bibitem{elsken2019nas}
T.~Elsken, J.~H. Metzen, and F.~Hutter,
Neural architecture search: a survey,
\textit{J.\ Mach.\ Learn.\ Res.}, 20 (2019), pp.~1--21.

\bibitem{liu2019darts}
H.~Liu, K.~Simonyan, and Y.~Yang,
DARTS: differentiable architecture search,
in Proc.\ ICLR 2019.

\bibitem{real2019aging}
E.~Real, A.~Aggarwal, Y.~Huang, and Q.~V. Le,
Regularized evolution for image classifier architecture search,
in Proc.\ AAAI 2019, vol.~33, pp.~4780--4789.

% ── General references ───────────────────────────────────────────
\bibitem{goldstein2002}
H.~Goldstein, C.~Poole, and J.~Safko,
\textit{Classical Mechanics}, 3rd ed.,
Addison-Wesley, Reading, MA, 2002.

\bibitem{lecun2015}
Y.~LeCun, Y.~Bengio, and G.~Hinton,
Deep learning,
\textit{Nature}, 521 (2015), pp.~436--444.

\bibitem{raiaan2024}
M.~A.~K. Raiaan \textit{et al.},
A systematic review of hyperparameter optimization techniques in CNNs,
\textit{Decision Anal.\ J.}, 11 (2024), p.~100470.

\bibitem{loshchilov2019adamw}
I.~Loshchilov and F.~Hutter,
Decoupled weight decay regularization,
in Proc.\ ICLR 2019.

\bibitem{paszke2019}
A.~Paszke \textit{et al.},
PyTorch: an imperative style, high-performance deep learning library,
in Adv.\ NeurIPS 32, 2019.

\bibitem{virtanen2020}
P.~Virtanen \textit{et al.},
SciPy 1.0: fundamental algorithms for scientific computing in Python,
\textit{Nat.\ Methods}, 17 (2020), pp.~261--272.

\bibitem{toth2020hamiltonian}
P.~Toth, D.~J. Rezende, A.~Jaegle, S.~Racani\`{e}re, A.~Botev,
and I.~Higgins,
Hamiltonian generative networks,
in Proc.\ ICLR 2020.

\bibitem{lindauer2020hpo}
M.~Lindauer and F.~Hutter,
Best practices for scientific research on neural architecture search,
\textit{J.\ Mach.\ Learn.\ Res.}, 21 (2020), pp.~1--18.

\bibitem{henon1964}
M.~H\'{e}non and C.~Heiles,
The applicability of the third integral of motion: some numerical
experiments,
\textit{Astron.\ J.}, 69 (1964), pp.~73--79.

\end{thebibliography}

\end{document}